% ==========================================
% LatexRefinement Step Output: step4-2026-10-06-tuesday-week4-algebra-1-offset-final.tex
% Model: gemini-3.8-flash
% Temperature: 0.35
% TopP: 0.9
% TopK: 10
% MaxOutputTokens: 65535
% ThinkingBudget: 4096
% ThinkingLevel: HIGH
% Prompt Tokens: 79’172
% Candidates Tokens: 38’383 (inkl. Thinking Tokens)
% Cached Tokens: 28’918
% Processed on: 2026-10-09 00:30:29
% ==========================================

% Primary Language: English
% End of the video: 01:30:14

\lecturechapter[4]{Tuesday}{06 Oct}{6 October 2026}{Polynomial Rings, Laurent Series, Quotient Rings, and the First Isomorphism Theorem}
% Old chapter title: Polynomial Rings, Laurent Series, and Ideals

\begin{overview}[Course Content Overview]
This lecture explores fundamental algebraic ring constructions, bridging polynomial extensions, localized series, and residue structures:
\begin{itemize}
    \item \textbf{Polynomial and Power Series Rings:} Multivariable polynomial rings, the binomial theorem in commutative rings, and the Frobenius identity in characteristic $p$.
    \item \textbf{Fields of Fractions and Laurent Series:} The field of rational functions $R(x)$ and formal Laurent series $R((x)) \cong \operatorname{Frac}(R[[x]])$.
    \item \textbf{Ideals and Quotient Rings:} Formal definition of ideals, the linear algebra analogy, and the construction of the quotient ring $R/I$ from scratch (highlighting how ideal absorption ensures well-defined coset multiplication).
    \item \textbf{The First Isomorphism Theorem for Rings:} Formulation and complete three-step proof of $R/\ker(f) \cong \operatorname{Im}(f)$.
    \item \textbf{Ideals Generated by Elements and Polynomial Quotients:} Contrasting the quotients $\mathbb{C}[x]/(x^2+1)$ (which possesses zero divisors) and $\mathbb{Q}[x]/(x^2+1) \cong \mathbb{Q}(i)$ (which forms the field of Gaussian rationals).
\end{itemize}
\end{overview}

\section{Review: Polynomial Rings and Power Series}

\begin{speech}[00:00:00 - 00:01:54]
Okay, so we start, and I wanted to start with some things from last week. If $R$ is a ring, we had defined this polynomial ring in one variable and also the power series ring. \inlinemetanote{writes $R \to R[x], R[[x]]$} And just some comments: the polynomial ring is a subring of the power series ring. \inlinemetanote{writes $R[x] \subset R[[x]]$} That's important to note, because every polynomial you can view as a finite power series.

So, this is one comment. The second comment --- and people were asking me after class --- is: what happens if I have two variables or three variables? There's two ways to approach this. One is: you could just define this... \inlinemetanote{writes $\to R[x,y] \to \text{defined}$} defined in the same way as we defined this, which is that it's the set of finite sums \inlinemetanote{writes $\{ \sum_{\text{finite}} a_{i,j} x^i y^j \mid a_{i,j} \in R \}$} $\sum a_{i,j} x^i y^j$, where these indices go over some finite set, and the coefficients are in $R$; and it's just the same definition as before. So, that's one way to do it. I hope that's clear --- it's just multiplying with more polynomials.

But there's a more formal way, which is that $R$ is a ring, so we've already defined this --- that's where we add one variable ---, and so, this is a ring, so we can also define this: this is a ring and we have one more variable. \inlinemetanote{writes $R \to R[x] \to (R[x])[y]$} And you can take this to be the definition of $R[x,y]$. \inlinemetanote{writes $= R[x,y]$}

So, we just add the variables one at a time. This is a polynomial ring over the polynomial ring over $R$, so it's another way. And, of course, if you have two different paths, then you have the burden of proving they're both the same, etc., and I leave this burden for you to carry, all right? \inlinemetanote{draws brace and writes \qt{same}} And these are the same; maybe I just say that: same.
\end{speech}

\begin{content}[Polynomial Rings in One and Several Variables]
Let $R$ be a commutative ring with identity. Recall the construction of the \newterm{polynomial ring} $R[x]$ and the \newterm{formal power series ring} $R[[x]]$. Every polynomial can be naturally identified with a formal power series having only finitely many non-zero coefficients, yielding the subring inclusion:
\[
R[x] \subseteq R[[x]].
\]

\begin{definition}[Polynomial Ring in Two Variables]\label[definition]{def:two-variable-polynomial-ring}
For two indeterminates $x$ and $y$, the polynomial ring $R[x, y]$ can be defined in two equivalent ways:
\begin{enumerate}[label=\textbf{(\arabic*)}]
    \item \textbf{Explicit Definition (Finite Linear Combinations):}
    \[
    R[x, y] := \left\{ \sum_{i, j \ge 0}^{\text{finite}} a_{i,j} x^i y^j \;\middle|\; a_{i,j} \in R \right\}.
    \]
    \item \textbf{Iterative Definition (Polynomials over Polynomials):}
    Viewing $R[x]$ as the coefficient ring, we form the univariate polynomial ring in $y$:
    \[
    R \longrightarrow R[x] \longrightarrow (R[x])[y] =: R[x, y].
    \]
\end{enumerate}
\end{definition}

\begin{steps}[Equivalence of Definitions]
Both constructions yield canonically isomorphic rings:
\[
(R[x])[y] \cong R[x, y] \cong (R[y])[x].
\]
In the iterative perspective, an element is a polynomial $P(y) = \sum_{j=0}^m p_j(x) y^j$, where each coefficient $p_j(x) = \sum_{i=0}^{k_j} a_{i,j} x^i \in R[x]$. Expanding via the distributive law expresses $P(y)$ uniquely as the finite double sum $\sum_{i, j} a_{i,j} x^i y^j$. By induction, this extends to $n$ variables:
\[
R[x_1, \dots, x_n] \cong (R[x_1, \dots, x_{n-1}])[x_n].
\]
\end{steps}
\end{content}

\subsection{The Binomial Formula in Commutative Rings}

\begin{speech}[00:01:54 - 00:03:47]
Okay, so then there's some comments. Essentially, well, something you should know is: if I have this polynomial ring in two variables and I take $x + y$, that's in this polynomial ring in two variables. And I can raise it to the $n$, and this is something you know --- this is this binomial expansion, or I don't know how it's called, binomial formula. \inlinemetanote{writes $R[x,y]$ and $(x+y)^n = \sum_{a=0}^n x^{n-a} y^a \binom{n}{a}$} This is equal to the sum --- so, from $a = 0$ to $n$ --- of $x^{n-a} y^a \binom{n}{a}$.

Okay? That's the binomial formula. Are you familiar with this binomial formula?

So, for example, $(x+y)^5$ is equal to $1 x^5 + 5 x^4 y + 10 x^3 y^2 + 10 x^2 y^3 + 5 x y^4 + y^5$. \inlinemetanote{writes $(x+y)^5 = x^5 + 5x^4y + 10x^3y^2 + 10x^2y^3 + 5xy^4 + y^5$} I guess I should have done a smaller one.

That's an example. And this holds in every ring, because these elements $x$ and $y$ are here, and we've already discussed that these numbers make sense in any ring --- they're just how you add up $1$. Is that clear? So, this formula holds in any ring. And to prove it, it's just the same as by induction: if you wanted to prove this formula in any ring, you just write it out as $(x+y)(x+y)\dots$ and use the distributive law, and you get $x^2$, and you regroup everything, and you get $xy + xy + y^2$, \inlinemetanote{writes $(x+y)^2 = (x+y)(x+y) = x^2 + xy + xy + y^2$} and that gives you the formula; and then you prove it by induction. Okay.
\end{speech}

\begin{content}[The Binomial Expansion]
Let $R$ be any commutative ring with identity $1_R$. For integers $k \in \mathbb{Z}$, the integer scalar $k \cdot 1_R$ is canonically defined by repeated addition or negation of $1_R$. Consequently, the binomial coefficients $\binom{n}{a} \in \mathbb{N}$ act canonically as coefficients in $R$.

\begin{theorem}[Binomial Formula in Commutative Rings]\label[theorem]{thm:binomial-formula}
In the polynomial ring $R[x, y]$, or for any two commuting elements $x, y \in R$ ($x y = y x$), the following identity holds for all $n \in \mathbb{N}$:
\begin{equation}\label{eq:binomial-formula}
(x + y)^n = \sum_{a=0}^n \binom{n}{a} x^{n-a} y^a.
\end{equation}
\end{theorem}

\begin{example}[Expansion for $n=5$ and $n=2$]\label[example]{ex:binomial-5}
For $n = 5$:
\[
(x + y)^5 = x^5 + 5 x^4 y + 10 x^3 y^2 + 10 x^2 y^3 + 5 x y^4 + y^5.
\]
For $n = 2$, distributivity and commutativity give:
\[
(x + y)^2 = (x + y)(x + y) = x^2 + x y + y x + y^2 = x^2 + 2 x y + y^2.
\]
\end{example}

\begin{steps}[Proof by Induction]
The general identity follows by induction on $n$:
\begin{align*}
(x + y)^{n+1} &= (x + y) \sum_{a=0}^n \binom{n}{a} x^{n-a} y^a \\
&= \sum_{a=0}^n \binom{n}{a} x^{n+1-a} y^a + \sum_{a=0}^n \binom{n}{a} x^{n-a} y^{a+1} \\
&= x^{n+1} + \sum_{k=1}^n \left[ \binom{n}{k} + \binom{n}{k-1} \right] x^{n+1-k} y^k + y^{n+1} \\
&= \sum_{k=0}^{n+1} \binom{n+1}{k} x^{n+1-k} y^k,
\end{align*}
using Pascal's identity $\binom{n}{k} + \binom{n}{k-1} = \binom{n+1}{k}$.
\end{steps}
\end{content}

\subsection{Characteristic \texorpdfstring{$p$}{p} and the Frobenius Identity}

\begin{speech}[00:03:47 - 00:05:43]
There's one comment here, which is kind of a fun comment. That holds in any ring, that binomial formula. But what happens if we take this ring $\mathbb{Z}/p\mathbb{Z}$, where $p$ is a prime? \inlinemetanote{writes $\mathbb{Z}/p\mathbb{Z}$, $p$ prime} And then in this ring, you could take this ring in $x, y$, \inlinemetanote{writes $(\mathbb{Z}/p\mathbb{Z})[x,y]$} and the fun thing to consider there is $(x+y)^p$. \inlinemetanote{writes $(x+y)^p$}

So, that's an example. And if I take this formula in $\mathbb{Z}/p\mathbb{Z}$ --- $\mathbb{Z}/5\mathbb{Z}$ in this case, where it's the same exponent as here --- something kind of nice happens: all these internal coefficients here, they all have a $5$ in them, so they all go away. And this is true just in general: in $\mathbb{Z}/p\mathbb{Z}$, if you take $(x+y)^p$, you get this formula that you secretly wanted anyway, which is this: \inlinemetanote{writes $(x+y)^p = x^p + y^p$} It's a simple formula.

And the reason is because for all of the binomial coefficients $\binom{p}{a}$, $p$ divides $\binom{p}{a}$ when $a$ goes from $1$ to $p-1$. \inlinemetanote{writes $p \mid \binom{p}{a}$ when $a \in \{1, \dots, p-1\}$} For all the internal coefficients, $p$ divides that. And why is that? It's because this binomial starts with a $p!\dots$ \inlinemetanote{writes $\binom{p}{a} = \frac{p!}{a!(p-a)!}$} and there are no $p$'s here to cancel the $p$ in the $p!$, because these are all lesser numbers. Is that clear?

So, you'll always have, if you look at the exponent being prime in this binomial theorem, all of the middle coefficients here will all have that prime in it. So, if I look at it in $\mathbb{Z}/p\mathbb{Z}$, I get this great formula... This is a well-known formula; it will come up. Okay, so that's an aside about the binomial formula.
\end{speech}

\begin{content}[The Freshman's Dream in Characteristic \texorpdfstring{$p$}{p}]
\begin{lemma}[Divisibility of Binomial Coefficients by a Prime]\label[lemma]{lem:prime-divides-binomial}
Let $p$ be a prime number. For every integer $a \in \{1, 2, \dots, p-1\}$, the prime $p$ divides the binomial coefficient $\binom{p}{a}$:
\begin{equation}\label{eq:prime-divides-binom}
p \mathrel{\Big|} \binom{p}{a}.
\end{equation}
\end{lemma}

\begin{proof}
Recall the algebraic formula:
\[
\binom{p}{a} = \frac{p!}{a! (p-a)!} \implies a! (p-a)! \binom{p}{a} = p! = p \cdot (p-1)!.
\]
Thus, $p$ divides the product $a! (p-a)! \binom{p}{a}$. Since $1 \le a \le p-1$, every factor occurring in the factorials $a! = 1 \cdot 2 \dots a$ and $(p-a)! = 1 \cdot 2 \dots (p-a)$ is strictly less than $p$. Because $p$ is prime, $\operatorname{gcd}(a! (p-a)!, p) = 1$. By Euclid's lemma, $p$ must divide the remaining factor $\binom{p}{a}$.
\end{proof}

\begin{proposition}[Frobenius Endomorphism Identity]\label[proposition]{prop:frobenius-identity}
Let $R$ be a commutative ring of characteristic $p$ (e.g., $R = \mathbb{Z}/p\mathbb{Z}$). In the polynomial ring $R[x, y]$, or for any commuting elements $x, y \in R$, we have:
\begin{equation}\label{eq:frobenius-freshmans-dream}
(x + y)^p = x^p + y^p.
\end{equation}
\end{proposition}

\begin{proof}
Expanding $(x + y)^p$ via \cref{thm:binomial-formula}:
\[
(x + y)^p = x^p + \sum_{a=1}^{p-1} \binom{p}{a} x^{p-a} y^a + y^p.
\]
By \cref{lem:prime-divides-binomial}, $\binom{p}{a} \equiv 0 \pmod p$ for all $1 \le a \le p-1$. In characteristic $p$, each intermediate term vanishes, leaving $(x + y)^p = x^p + y^p$.
\end{proof}
\end{content}

\begin{hidden}
% timestamp: 00:05:43
% topic: Review of multivariable polynomials, binomial theorem, and Frobenius identity in characteristic p.
% board_state: def:two-variable-polynomial-ring, thm:binomial-formula, lem:prime-divides-binomial, prop:frobenius-identity
% next_goal: Discuss integral domains for R[x] and R[[x]], and fields of fractions.
% open_loops: none
\end{hidden}

\section{Integral Domains, Fields of Fractions, and Laurent Series}

\subsection{Polynomial Rings and Power Series Rings as Integral Domains}

\begin{speech}[00:05:43 - 00:06:38]
And then there were some definitions: that if I take $R$ now to be an integral domain, \inlinemetanote{writes $R$ is domain} we had defined here $R[x]$ --- that's polynomials in one variable ---; this is also an integral domain. \inlinemetanote{writes $\implies R[x]$ also int.\ domain} We discussed that: you use the highest term, and if $R$ is an integral domain, when you multiply polynomials, the highest term lives, so it stays an integral domain.

What I didn't say is: this is also an integral domain. \inlinemetanote{writes $R[[x]]$ also int.\ domain} And one needs a different argument here, because power series don't have highest terms. And so, the answer is... well, what do you think the answer is? You can't use the highest term here, so what should you use?
\end{speech}

\begin{student}[Student Answer]
Lowest term.
\end{student}

\begin{speech}[00:06:40 - 00:07:07]
Lowest term, exactly! You could have used the lowest term here, but psychologically, with polynomials, you go to the highest term. But it turns out lowest term is a better idea here, because a power series does have a lowest term. And when you multiply the two lowest terms of power series, it's the lowest term of the product. So, it's the same argument here: it's just using the lowest term. Which is, you know, psychologically a little strange --- you could have used the lowest term here, but you didn't, right? Okay, so this is also an integral domain.
\end{speech}

\begin{content}[Integral Domain Properties: Degrees vs. Orders]
\begin{theorem}[Preservation of the Integral Domain Property]\label[theorem]{thm:domain-poly-power-series}
If $R$ is an integral domain, then:
\begin{enumerate}[label=\textbf{(\arabic*)}]
    \item The polynomial ring $R[x]$ is an integral domain.
    \item The formal power series ring $R[[x]]$ is an integral domain.
\end{enumerate}
\end{theorem}

\begin{proof}
Let $f, g$ be non-zero elements in the respective ring.
\begin{enumerate}[label=\textbf{(\arabic*)}]
    \item \textbf{For $R[x]$ (Highest Degree Term):}
    Let $f = a_n x^n + \dots + a_0$ and $g = b_m x^m + \dots + b_0$ with $a_n \neq 0$ and $b_m \neq 0$. The leading term of the product $f g$ is:
    \[
    a_n b_m x^{n+m}.
    \]
    Since $R$ is an integral domain, $a_n b_m \neq 0$. Thus, $f g \neq 0$ and $\deg(f g) = \deg(f) + \deg(g)$.
    \item \textbf{For $R[[x]]$ (Lowest Degree / Order Term):}
    A formal power series does not have a maximal degree, but every non-zero series $f = \sum_{k=0}^\infty a_k x^k$ has a well-defined \newterm{order} (minimal non-zero index):
    \[
    \operatorname{ord}(f) := \min \{ k \in \mathbb{N} \mid a_k \neq 0 \}.
    \]
    We write $f = x^n (a_n + a_{n+1} x + \dots)$ with $a_n \neq 0$, and $g = x^m (b_m + b_{m+1} x + \dots)$ with $b_m \neq 0$. The lowest non-vanishing term of the Cauchy product is:
    \[
    a_n b_m x^{n+m}.
    \]
    Since $R$ is an integral domain, $a_n b_m \neq 0$. Consequently, $f g \neq 0$ with $\operatorname{ord}(f g) = \operatorname{ord}(f) + \operatorname{ord}(g)$.
\end{enumerate}
\end{proof}
\end{content}

\subsection{Fields of Fractions: Rational Functions and Laurent Series}

\begin{speech}[00:07:07 - 00:08:35]
And the point I wanted to make here is that the main construction last week was the field of fractions. \inlinemetanote{draws arrow and writes \qt{field of fractions}} We'll have a different construction this week, but that was the field of fractions; and the important thing about a field of fractions is: you can only build that on an integral domain.

So, if I have an integral domain, I could take the field of fractions of that integral domain --- that was the construction. But I could also take the field of fractions of polynomials, because that's an integral domain; and that is called... well, the symbol for that is $R$ with round parentheses --- that's round brackets ---; that's the name of the field of fractions of the polynomial ring when $R$ is an integral domain. And it's called the field of rational functions. \inlinemetanote{writes \qt{field of rational functions}}

So, rational functions with coefficients in $R$. This is polynomials with coefficients in $R$; that's the field of rational functions. And you know how to think about this because we've thought about fields of fractions: an element here in rational functions is a polynomial divided by another polynomial. \inlinemetanote{writes $R(x) \ni \frac{f(x)}{g(x)}$} So, both of these, $f$ and $g$, are polynomials, \inlinemetanote{writes $f, g \in R[x]$} and, of course, $g$ is not allowed to be zero. \inlinemetanote{writes $g \neq 0$} So, that's how to think about an element here, and that's...
\end{speech}

\begin{content}[The Field of Rational Functions]
Recall that for any integral domain $D$, its \newterm{field of fractions} $\operatorname{Frac}(D)$ consists of equivalence classes of formal quotients $\frac{a}{b}$ with $a, b \in D$ and $b \neq 0$.

\begin{definition}[Field of Rational Functions $R(x)$]\label[definition]{def:rational-functions}
Let $R$ be an integral domain. The field of fractions of the polynomial ring $R[x]$ is denoted by round parentheses $R(x)$ and is called the \newterm{field of rational functions}:
\begin{equation}\label{eq:field-rational-functions}
R(x) := \operatorname{Frac}(R[x]) = \left\{ \frac{f(x)}{g(x)} \;\middle|\; f(x), g(x) \in R[x], \ g(x) \neq 0 \right\}.
\end{equation}
\end{definition}
\end{content}

\begin{speech}[00:08:35 - 00:10:20]
Now, power series are also integral domains, so we can again consider the field of fractions, and the symbol for that is two parentheses. \inlinemetanote{writes $R((x))$} All right, so this is kind of universal math notation: you have an integral domain, you have the polynomial ring, power series, field of fractions, and this one has a special name. Does anyone know what the name of this thing is? It's called the field of Laurent series. \inlinemetanote{writes \qt{Field of Laurent Series}} That's the name of this.

Actually, I'm having some doubts about how to spell Laurent. Is it with an \qt{a}? No, I think it's with an \qt{e}. I'm not $100\%$ sure about that.

So, this is the field of Laurent series, and this will come up in your complex analysis classes for sure, but with a $\mathbb{C}$ here. And usually the argument is called $z$ in complex analysis.

So, the field of Laurent series... you could say: we've already constructed the field of fractions, so why do I have to talk about it here? You could say this field of Laurent series we think about in the same way, as $f(x)$ over $g(x)$, where $f$ and $g$ are power series. \inlinemetanote{writes $R((x)) \ni \frac{f(x)}{g(x)}$} And this would be absolutely correct. \inlinemetanote{writes $f, g \in R[[x]], g \neq 0$} That would be correct, with $g$ not equal to zero.

It's just there's a nicer way to think about it, and that has to do with the last thing I said --- not yesterday, last week. I think it's with \qt{e}.
\end{speech}

\begin{content}[The Field of Fractions of Formal Power Series]
\begin{definition}[Field of Fractions $R((x))$]\label[definition]{def:frac-power-series}
Let $R$ be an integral domain. Since $R[[x]]$ is an integral domain (\cref{thm:domain-poly-power-series}), we define its field of fractions, denoted with double round parentheses $R((x))$:
\begin{equation}\label{eq:frac-power-series}
R((x)) := \operatorname{Frac}(R[[x]]) = \left\{ \frac{f(x)}{g(x)} \;\middle|\; f(x), g(x) \in R[[x]], \ g(x) \neq 0 \right\}.
\end{equation}
This field is known algebraically as the \newterm{field of formal Laurent series}.
\end{definition}
\end{content}

\begin{note}[Board Cleaning Transition]
\inlinemetanote{The lecturer pauses to erase the middle and right boards, asking the class if they have already encountered Laurent series in complex analysis.}
\end{note}

\subsection{Formal Laurent Series with Finitely Many Negative Powers}

\begin{speech}[00:11:02 - 00:13:20]
You won't go wrong if you think about it this way: that's the definition, that's how we constructed it. But there's another way, which is why it's called the field of Laurent series. So, what is a Laurent series? \inlinemetanote{writes \qt{Def: Laurent series in $x$}}

A Laurent series is a power series where you are allowed to have finitely many negative coefficients --- that's the definition. So, I can say this: a Laurent series in a variable $x$ is a power series with --- and this is the crucial thing --- finitely many negative powers. \inlinemetanote{writes \qt{is a power series with finitely many negative powers}} You can have as many positive-degree coefficients as you want, it's a power series, but only finitely many negative powers.

So, that's the crucial thing for a Laurent series. What is a power series? A power series might be like the power series for $e^x$; we've seen these anyway: \inlinemetanote{writes $1 + x + \frac{x^2}{2!} + \frac{x^3}{3!} + \dots \in R[[x]]$} That's a power series. These are elements of the power series ring.

In a Laurent series, you're allowed to have some negative powers, so you could have $x^{-12} + 3x^{-2} + x^{-1}$, then the constant term... \inlinemetanote{writes $x^{-12} + 3x^{-2} + x^{-1} + 17 - x + x^{15} + \dots$} I got too excited about the constant term: $17 - x + x^{15} + \dots$ That is a Laurent series.

Have you seen these objects in mathematics? Okay, well, maybe it's the first time, but you can imagine: you just take power series, and you're allowed to have negative terms, but only finitely many. And it's extremely important that you have only finitely many. You'll find that if you have only finitely many, you can add them, and when you add them you'll still only have finitely many, because each of the two only had finitely many. When you multiply them, the lowest term now... if I multiply this with something with lowest term $-7$, its lowest term will be $x^{-19}$, so it will still have only finitely many. Is that clear?
\end{speech}

\begin{content}[Formal Definition and Arithmetic of Laurent Series]
\begin{definition}[Formal Laurent Series]\label[definition]{def:laurent-series}
A \newterm{formal Laurent series} in the indeterminate $x$ over a ring $R$ is a formal sum of the form:
\begin{equation}\label{eq:laurent-series-def}
f(x) = \sum_{n = -N}^\infty a_n x^n = a_{-N} x^{-N} + a_{-N+1} x^{-N+1} + \dots + a_0 + a_1 x + a_2 x^2 + \dots
\end{equation}
where $N \in \mathbb{Z}$ and $a_n \in R$ for all $n \ge -N$. That is, $f(x)$ contains at most \textbf{finitely many negative powers} of $x$.
\end{definition}

\begin{example}[Power Series vs. Laurent Series]\label[example]{ex:laurent-series}
\begin{itemize}
    \item \textbf{Power Series:} Elements of $R[[x]]$ contain only non-negative powers:
    \[
    1 + x + \frac{x^2}{2!} + \frac{x^3}{3!} + \dots \in R[[x]].
    \]
    \item \textbf{Laurent Series:} Elements can have negative exponents up to some lowest power:
    \[
    x^{-12} + 3 x^{-2} + x^{-1} + 17 - x + x^{15} + \dots \in R((x)).
    \]
\end{itemize}
\end{example}
\end{content}

\begin{speech}[00:13:20 - 00:15:20]
And here's the great thing: if you're non-zero here, you can divide. That's why the Laurent series turns out to be a field. And this is something you should think about --- we did discuss this last time.

So, if you have a power series... maybe I just say this, because this is maybe now departing a little bit from the course, but it is the case that if I take this field of Laurent series... \inlinemetanote{writes \qt{field of Laurent series $\cong R((x))$}} this field of Laurent series is isomorphic to $R((x))$, the field of fractions of the power series. And what you have to show there is that if you have Laurent series, you can divide; that's the interesting thing. And that has to do with the last thing I said last week, which is that it's pretty easy to divide power series: all you have to do is find the first non-zero coefficient, then use this rule I explained last week, and then you can play with this power to take it out. So, I leave this for you.

I'll say one thing: why don't we allow power series with infinitely many negative terms? Because if you try to multiply them, then you have to do infinite sums. And we don't do infinite sums in algebra; we only do finite sums. Because if you do infinite sums, you need convergence, and you can try to do that --- that's an interesting subject ---, but in this class we're not doing infinite sums. And if you have finitely many negative terms, when you multiply, you see you'll never have to do infinite sums; it's just the same.

So, anyway, that's just this comment here: this field of Laurent series, it turns out, is just a nice way of thinking about this field of fractions for power series. And you can imagine it's nicer to think about this kind of object than the ratio. And this second way of looking at things is not available for polynomials: there's only one way to think about it there, as ratios of polynomials. Okay.
\end{speech}

\begin{content}[Isomorphism between Laurent Series and the Field of Fractions]
\begin{theorem}[Algebraic Characterization of $R((x))$]\label[theorem]{thm:laurent-field-isomorphism}
When $K$ is a field, the set of formal Laurent series with finitely many negative powers forms a field under formal addition and multiplication, and is naturally isomorphic to the field of fractions of $K[[x]]$:
\begin{equation}\label{eq:laurent-isomorphism}
\operatorname{Frac}(K[[x]]) \cong K((x)) = K[[x]][x^{-1}].
\end{equation}
\end{theorem}

\begin{steps}[Why Invertibility Holds and Why Infinitely Many Negative Terms Fail]
\begin{enumerate}[label=\textbf{(\arabic*)}]
    \item \textbf{Invertibility in $K((x))$:}
    Any non-zero Laurent series can be factored as $f(x) = x^k u(x)$, where $k \in \mathbb{Z}$ is the minimal index with $a_k \neq 0$, and $u(x) = a_k + a_{k+1} x + \dots \in K[[x]]$ with $a_k \neq 0$. Since $K$ is a field, the leading coefficient $a_k$ is a unit in $K$, so $u(x)$ is an invertible unit in the power series ring $K[[x]]$. Its multiplicative inverse is:
    \[
    f(x)^{-1} = x^{-k} u(x)^{-1} \in K((x)).
    \]
    \item \textbf{Failure of Infinitely Many Negative Terms:}
    If series with infinitely many negative powers were permitted ($\sum_{n=-\infty}^\infty a_n x^n$), the coefficient of $x^k$ in the product of two such series would require evaluating the infinite sum:
    \[
    c_k = \sum_{j=-\infty}^\infty a_j b_{k-j}.
    \]
    Without topological convergence, such an infinite sum is undefined in pure algebra. Restricting to finitely many negative powers ensures that for each $k$, only finitely many non-zero products $a_j b_{k-j}$ occur, keeping all algebraic operations purely finite.
\end{enumerate}
\end{steps}
\end{content}

\begin{hidden}
% timestamp: 00:15:20
% topic: Laurent series as field of fractions R((x)); introduction of ideals and quotient rings.
% board_state: def:rational-functions, def:frac-power-series, def:laurent-series, thm:laurent-field-isomorphism
% next_goal: Formal definition of ideals and ring homomorphisms.
% open_loops: none
\end{hidden}

\section{Ideals and Ring Homomorphisms}

\subsection{Definition and Axioms of an Ideal}

\begin{speech}[00:15:33 - 00:16:29]
Okay, so the main topic this week is about ideals and quotient rings, and the main construction is that of a quotient ring.

If you like polynomial rings, the class will talk about results of polynomial rings, but mostly polynomial rings in one variable. The subject of polynomial rings in several variables is very important in mathematics, and there's a whole class about it, more or less: the class is called commutative algebra.
\end{speech}

\begin{note}[Board Cleaning and Transition]
\inlinemetanote{The lecturer finishes erasing the boards to begin the new chapter on ideals and quotient rings.}
\end{note}

\begin{speech}[00:16:30 - 00:18:33]
Okay, so now we finish these constructions and we start something new. So, $R$ is our ring, \inlinemetanote{writes $R$ ring} and the first thing, which we already defined: $I$ is an ideal. \inlinemetanote{writes $I \subset R$ is an ideal} We already defined this, but I'll just review it briefly, because now it's crucial for this week.

$I$ is an ideal if it satisfies that if I take $(I, +, 0)$, this is a subgroup of the additive part of the ring. \inlinemetanote{writes $\bullet \ (I, +, 0) \subset (R, +, 0)$ subgroup} And the second is: if I take anybody in the ring and anybody in the ideal, then their product lies in the ideal. \inlinemetanote{writes $\bullet \ \forall r \in R, \forall x \in I \text{ then } r \cdot x \in I$} This says that the ideal is an additive subgroup under addition, but there's another condition. The ideal has two constraints: one is the addition constraint, and the other is the multiplication constraint. They're just different, because they're talking about two different operations.

And our example we had seen --- we did some examples anyway ---, but a place where we know we always get ideals: if I take a homomorphism... \inlinemetanote{writes $f \colon R \to S$} Maybe I should put it here: so, $f$ is a homomorphism of rings, \inlinemetanote{writes $f$ is a hom of rings} then whenever you have a homomorphism of rings, you get an ideal for free. Sorry. \inlinemetanote{drops chalk}

And we discussed this also before: then the kernel of $f$ --- which is defined to be all the elements in the ring such that $f(a)$ is equal to $0$ \inlinemetanote{writes $\ker(f) := \{ a \in R \mid f(a) = 0 \}$} --- is an ideal. \inlinemetanote{writes \qt{is an ideal}} We proved this in one of the early weeks.
\end{speech}

\begin{content}[Definition of Ideals and the Kernel of a Homomorphism]
Let $R$ be a commutative ring with identity.

\begin{definition}[Ideal]\label[definition]{def:ideal}
A subset $I \subseteq R$ is called an \newterm{ideal} of $R$ (denoted $I \trianglelefteq R$) if it satisfies the following two conditions:
\begin{enumerate}[label=\textbf{(\arabic*)}]
    \item \textbf{Additive Subgroup Condition:}
    $(I, +, 0)$ is a subgroup of the additive group $(R, +, 0)$. That is, $0 \in I$, and for all $x, y \in I$, $x - y \in I$.
    \item \textbf{Absorption / Multiplication Constraint:}
    $I$ absorbs multiplication by arbitrary elements of the ambient ring $R$:
    \[
    \forall r \in R, \quad \forall x \in I \implies r \cdot x \in I.
    \]
\end{enumerate}
\end{definition}

\begin{proposition}[Kernel of a Ring Homomorphism as an Ideal]\label[proposition]{prop:kernel-is-ideal}
Let $f \colon R \to S$ be a homomorphism of commutative rings. The \newterm{kernel} of $f$, defined by:
\begin{equation}\label{eq:kernel-definition}
\ker(f) := \{ a \in R \mid f(a) = 0_S \},
\end{equation}
is an ideal of $R$.
\end{proposition}

\begin{proof}
We verify the two defining conditions of \cref{def:ideal}:
\begin{enumerate}[label=\textbf{(\arabic*)}]
    \item \textbf{Additive Subgroup:} Since $f(0_R) = 0_S$, $0_R \in \ker(f)$. For $x, y \in \ker(f)$, additivity of $f$ yields $f(x - y) = f(x) - f(y) = 0_S - 0_S = 0_S$, so $x - y \in \ker(f)$.
    \item \textbf{Multiplicative Absorption:} For any $r \in R$ and $x \in \ker(f)$, multiplicativity gives:
    \[
    f(r \cdot x) = f(r) \cdot f(x) = f(r) \cdot 0_S = 0_S.
    \]
    Thus, $r \cdot x \in \ker(f)$.
\end{enumerate}
\end{proof}
\end{content}

\subsection{The Linear Algebra Analogy: Kernels, Images, and Isomorphisms}

\begin{speech}[00:18:33 - 00:20:20]
So, in general, when you study algebra, even in linear algebra --- I remind you in the study of linear algebra, this is a reminder: \inlinemetanote{writes \qt{Reminder:}}

If you have a vector space $V$ and you have a linear transformation, let's call it $\phi$. \inlinemetanote{writes $\phi \colon V \to W$} So, $V$ and $W$ are vector spaces \inlinemetanote{writes $V, W$ vector spaces} over some field --- it's not relevant for this example, say the complex numbers ---, and $\phi$ is a linear transformation. \inlinemetanote{writes $\phi$ linear transformation} I don't know what you called homomorphisms of vector spaces: linear transformations.

In the abstract study of linear algebra, when you have such a linear transformation, you automatically get two other vector spaces, so to speak for free. The two vector spaces are the kernel \inlinemetanote{writes $\ker(\phi) \subseteq V$} --- this is a subspace of $V$ ---, and the image, \inlinemetanote{writes $\operatorname{Im}(\phi) \subseteq W$} which is a subspace of $W$. And there is a theorem that links this data. If you remember, it says that if I take $V$ and I quotient by the kernel --- I hope you remember this, it's kind of crucial here ---, if you take $V$ and you quotient by the kernel, this is isomorphic as vector spaces to the image. \inlinemetanote{writes $V / \ker(\phi) \cong \operatorname{Im}(\phi)$}

All right, so this is kind of a model. Did you learn it in these terms? I've never taught linear algebra here. Okay. So, we would like to think a little bit about this kind of construction for rings, and we'll see a lot of it to start.
\end{speech}

\begin{overview}[Linear Algebra Analogy: The First Isomorphism Theorem]
For a linear transformation $\phi \colon V \to W$ between vector spaces over a field $\mathbb{F}$:
\begin{itemize}
    \item \textbf{Kernel:} $\ker(\phi) = \{ v \in V \mid \phi(v) = 0 \} \le V$ is a vector subspace of $V$.
    \item \textbf{Image:} $\operatorname{Im}(\phi) = \{ \phi(v) \mid v \in V \} \le W$ is a vector subspace of $W$.
    \item \textbf{First Isomorphism Theorem for Vector Spaces:}
    \[
    \frac{V}{\ker(\phi)} \cong \operatorname{Im}(\phi).
    \]
\end{itemize}
In ring theory, we seek an identical structural correspondence for a ring homomorphism $f \colon R \to S$.
\end{overview}

\begin{speech}[00:20:20 - 00:21:11]
Yeah, I haven't taught linear algebra for... I don't know, it's possible not in this century. But it's a great class, linear algebra. I used to tell students it's like the most important math class, linear algebra, because humans are actually pretty good at solving linear equations and pretty terrible at solving any other kind of equation.

That's not true, I did teach it. \inlinemetanote{cleans and wipes the lower board}
\end{speech}

\subsection{Ring Homomorphisms vs. Linear Transformations}

\begin{speech}[00:21:11 - 00:23:08]
Okay, so if I have this ring homomorphism... \inlinemetanote{writes $f \colon R \to S$} So, these are both rings now --- I try to make the parallel here: $R$ and $S$ are rings, \inlinemetanote{writes $R, S$ are rings} and I remind you they're commutative rings with unit, and $f$ always takes unit: $f$ takes $1$ to $1$, that's the definition of a ring homomorphism. And $f$ is a homomorphism, \inlinemetanote{writes $f$ hom} so we just want to complete that board.

Now, in this situation, we get two objects for free --- or, I don't know, \qt{free} might be the wrong word, but anyway, this gives us two objects. One is the kernel inside $R$, but this is not... there, everything was vector spaces, but this is not a ring now: this is an ideal. \inlinemetanote{writes $\ker(f) \subseteq R$ ideal} But we also get the image of $f$ inside $S$, and that is a subring. \inlinemetanote{writes $\operatorname{Im}(f) \subseteq S$ subring} We saw that the image was a subring.

So, that's kind of the parallel there. And then we're going to work on this quotient business. But you see, the matter is slightly... in this vector space world, everything is a vector space. But in the ring world, sometimes you get rings and sometimes you get ideals, and they're slightly different objects. So, we have to come to terms with that.

Let me just say some basic statements: \inlinemetanote{writes \qt{statements:}} One is that $f$ is surjective if and only if this subring is the whole ring. \inlinemetanote{writes \textbf{(1)} $f$ is surjective $\iff \operatorname{Im}(f) = S$} That's almost like a tautology, right? That's what surjective means: you get everybody here, and if the image is everybody, it's surjective. That's essentially a tautology.
\end{speech}

\begin{content}[The Morphism Data for Rings: Ideals vs. Subrings]
When comparing linear transformations $\phi \colon V \to W$ and ring homomorphisms $f \colon R \to S$, a crucial categorical divergence arises:
\begin{center}
\begin{tabular}{l l l}
    \toprule
    \textbf{Structure} & \textbf{Linear Algebra ($\phi \colon V \to W$)} & \textbf{Ring Theory ($f \colon R \to S$)} \\
    \midrule
    \textbf{Domain Sub-object} & $\ker(\phi) \le V$ (Vector Subspace) & $\ker(f) \trianglelefteq R$ (\textbf{Ideal}, not a subring with $1_R$) \\
    \textbf{Codomain Sub-object} & $\operatorname{Im}(\phi) \le W$ (Vector Subspace) & $\operatorname{Im}(f) \le S$ (\textbf{Subring}, contains $1_S = f(1_R)$) \\
    \bottomrule
\end{tabular}
\end{center}

\begin{proposition}[Surjectivity Criterion]\label[proposition]{prop:surjective-image}
Let $f \colon R \to S$ be a ring homomorphism. Then $f$ is surjective if and only if:
\begin{equation}\label{eq:surjective-criterion}
\operatorname{Im}(f) = S.
\end{equation}
\end{proposition}
\end{content}

\begin{hidden}
% timestamp: 00:23:08
% topic: Parallel between linear transformations and ring homomorphisms; properties of kernel and image.
% board_state: def:ideal, prop:kernel-is-ideal, prop:surjective-image
% next_goal: Characterization of injectivity via the zero ideal (statement 2).
% open_loops: none
\end{hidden}

\subsection{Injectivity and the Zero Ideal}

\begin{speech}[00:23:08 - 00:24:24]
The second statement --- and that's also true there --- is: $f$ is injective if and only if... \inlinemetanote{writes \textbf{(2)} $f$ is injective $\iff \ker(f) = (0)$}

So, what's the statement going to be? It's going to be about the kernel: the kernel of $f$ is the zero ideal. This is the standard terminology for the zero ideal, the standard notation for the zero ideal. The zero ideal is the ideal that has $0$ and nothing else. If you look at the definition of an ideal, the ideal that has the single element $0$... This is going to be used a lot, so it should be completely clear: \inlinemetanote{writes \qt{Zero ideal: $(0)$, $0$, $\{0\}$}}

The zero ideal: sometimes this is denoted like this, $(0)$; sometimes it's just denoted $0$; but the zero ideal is the ideal that contains one element: $\{0\}$. \inlinemetanote{writes $\{0\}$} It's definitely a subgroup, because it's just zero: zero, when you add it, stays zero, it has inverses, there's nothing wrong with zero. And the second condition is: it's closed under multiplication, because if you take zero in the ideal and multiply by anybody, you stay zero. So, that's the zero ideal.

And this last statement says: $f$ is injective if and only if the kernel is the zero ideal.
\end{speech}

\begin{content}[Injectivity and the Zero Ideal]
\begin{definition}[The Zero Ideal]\label[definition]{def:zero-ideal}
In any ring $R$, the singleton set consisting solely of the additive neutral element:
\begin{equation}\label{eq:zero-ideal-def}
(0) = \{0\}
\end{equation}
is an ideal of $R$, called the \newterm{zero ideal}.
\end{definition}

\begin{proposition}[Injectivity and the Kernel]\label[proposition]{prop:injective-kernel}
Let $f \colon R \to S$ be a ring homomorphism. Then $f$ is injective if and only if its kernel is trivial:
\begin{equation}\label{eq:injective-kernel-criterion}
f \text{ is injective} \iff \ker(f) = (0).
\end{equation}
\end{proposition}
\end{content}

\begin{proofWrapper}[Proof of \cref{prop:injective-kernel}]

\begin{speech}[00:24:24 - 00:25:39]
So, let's try to see what the argument for that is --- that's almost a tautology. Proof of \textbf{(2)}: \inlinemetanote{writes \qt{Proof of (2):}}

Suppose... we have to do two things. Suppose $f$ is injective. \inlinemetanote{writes $(\implies)$ Suppose $f$ is injective} Then let $a \in R$ be non-zero. \inlinemetanote{writes let $a \in R, a \neq 0$} Then $f(a)$ is not equal to $f(0)$, \inlinemetanote{writes then $f(a) \neq f(0)$} because injective means if you take different elements, they go to different places. And $f(0) = 0$, \inlinemetanote{writes $= 0$} so that means $f(a)$ is not zero, which implies $a$ is not in the kernel of $f$. \inlinemetanote{writes $\implies a \notin \ker(f)$}

So, if $f$ is injective, nobody who is not zero is in the kernel of $f$, so this implies the kernel of $f$ equals $(0)$. \inlinemetanote{writes $\implies \ker(f) = (0)$} That's one half of it: I assume it's injective, and I deduce the kernel of $f$ is zero --- it can only have one element in it.
\end{speech}

\begin{content}[Forward Direction: Injectivity Implies Trivial Kernel]
\begin{proof}[Forward Implication ($\implies$)]
Suppose $f$ is injective. Let $a \in R$ with $a \neq 0$.
Since $f$ is injective and $a \neq 0$:
\[
f(a) \neq f(0).
\]
Since $f$ is a ring homomorphism, $f(0) = 0_S$. Thus:
\[
f(a) \neq 0_S \implies a \notin \ker(f).
\]
Consequently, no non-zero element of $R$ belongs to $\ker(f)$. Since $0 \in \ker(f)$ always holds, we conclude:
\[
\ker(f) = (0).
\]
\end{proof}
\end{content}

\begin{speech}[00:25:39 - 00:26:39]
The other way around is: suppose the kernel of $f$ is $(0)$ --- \inlinemetanote{writes $(\impliedby)$ Suppose $\ker(f) = (0)$} and as I said, it's standard to use both notations, depending on what you want to emphasize.

Suppose the kernel of $f$ is zero, and let $f(a) = f(b)$. \inlinemetanote{writes then let $f(a) = f(b)$} We're assuming it's zero, and then we pick two elements $a, b$ in the ring that go to the same place. This implies, by the homomorphism property, $f(a) - f(b) = f(a - b) = 0$, because I'm assuming they're the same. \inlinemetanote{writes $\implies f(a) - f(b) = f(a-b) = 0$}

So, this implies that $a - b$ is in the kernel; \inlinemetanote{writes $\implies a - b \in \ker(f)$} but we're assuming the kernel is zero, so this implies $a - b = 0$, which implies $a = b$. \inlinemetanote{writes $\implies a - b = 0 \implies a = b$}

That's the second implication. And actually, if you look at this proof, it's the same proof as you had in linear algebra; there's actually no difference.
\end{speech}

\begin{content}[Reverse Direction: Trivial Kernel Implies Injectivity]
\begin{proof}[Reverse Implication ($\impliedby$)]
Suppose $\ker(f) = (0)$. Let $a, b \in R$ such that $f(a) = f(b)$.
Using the additive homomorphism property of $f$:
\[
f(a) - f(b) = 0_S \implies f(a - b) = 0_S.
\]
By definition of the kernel, this implies:
\[
a - b \in \ker(f).
\]
Since $\ker(f) = (0) = \{0\}$ by hypothesis, it follows that:
\[
a - b = 0_R \implies a = b.
\]
Thus, $f$ is injective, completing the proof of \cref{prop:injective-kernel}.
\end{proof}
\end{content}

\end{proofWrapper}

\begin{hidden}
% timestamp: 00:26:39
% topic: Equivalence of injectivity and trivial kernel; parallel to linear algebra.
% board_state: prop:surjective-image, def:zero-ideal, prop:injective-kernel
% next_goal: Transition to the quotient ring construction R/I.
% open_loops: none
\end{hidden}

\begin{speech}[00:26:39 - 00:28:05]
For this linear algebra picture, we've recovered all the pieces except this one, \inlinemetanote{points to $V/\ker(\phi) \cong \operatorname{Im}(\phi)$ on the upper right board} which says that $V$ mod the kernel is the image, and we now want to recover that. And for that, we need a construction. There was a construction already in linear algebra you needed, which was the notion of a quotient vector space; if you remember that, it's going to be very helpful.

The idea of a math education is that, actually, you're not allowed to forget anything you've learned. \inlinemetanote{erases the board} That's one theory. And the other theory is that education is what's left after you've forgotten everything you've learned. You can decide which theory of education you'd like to subscribe to. \inlinemetanote{cleans the lower board completely with a sponge}
\end{speech}

\begin{note}[Board Cleaning and Pedagogical Aside]
The lecturer erases the lower blackboard to prepare for the construction of quotient rings, sharing an aside on two contrasting theories of education.
\end{note}

\section{The Construction of Quotient Rings}

\begin{speech}[00:28:05 - 00:29:35]
Okay, so the task at hand is to understand this quotient. And now we do this not just for the kernel --- we'll get back to the kernel ---, but for any ideal. So, we let $I \subseteq R$ be an ideal. \inlinemetanote{writes Let $I \subseteq R$ be an ideal} And now this is the main construction for this week: we will construct the quotient ring. \inlinemetanote{writes We will construct the quotient ring} That's the name: quotient ring. And the symbol is $R/I$. \inlinemetanote{writes $R/I$} That's the symbol.

We have to construct it, and the parallel to linear algebra, of course, is $R$ mod the kernel --- that's the quotient ring. We'll get back to that in a moment. So, we forget about this linear algebra analogy for now, and we go for this construction: $R$ is a ring, $I$ is an ideal, and we need to construct the quotient.

We need to construct this object as a ring. That means, as usual, it's a multi-step construction: we have to first construct it as a set; then we have to find zero in it; we have to find one in it, \inlinemetanote{writes $(R/I, +, \cdot, 0, 1)$} zero and one; and we have to define addition and multiplication. You have to construct the whole thing --- that's just what you have to do; there's no other way around it.
\end{speech}

\begin{content}[The Quotient Ring Construction: Roadmap]
Let $R$ be a commutative ring with identity $1$, and let $I \trianglelefteq R$ be an ideal. Our objective is to construct a new ring:
\begin{equation}\label{eq:quotient-ring-symbol}
R/I = (R/I, +, \cdot, 0_{R/I}, 1_{R/I}),
\end{equation}
called the \newterm{quotient ring} (or residue class ring) of $R$ modulo $I$.

\begin{steps}[The Five Stages of Constructing a Ring from Scratch]
To endow the quotient $R/I$ with a complete ring structure, we must rigorously execute the following stages:
\begin{enumerate}[label=\textbf{Step \arabic*:}]
    \setcounter{enumi}{0}
    \item \textbf{Underlying Set:} Define an equivalence relation on $R$ and take $R/I$ to be the set of its equivalence classes.
    \item \textbf{Distinguished Elements:} Identify the additive zero $0_{R/I}$ and multiplicative identity $1_{R/I}$.
    \item \textbf{Addition Operation:} Define binary addition $+$ on equivalence classes and prove it is \emph{well-defined}.
    \item \textbf{Multiplication Operation:} Define binary multiplication $\cdot$ on equivalence classes and prove it is \emph{well-defined} using the ideal absorption property.
    \item \textbf{Ring Axioms Verification:} Verify associativity, commutativity, distributivity, and existence of inverses.
\end{enumerate}
\end{steps}
\end{content}

\subsection{Step 1: The Underlying Set and Equivalence Relation Modulo an Ideal}

\begin{speech}[00:29:35 - 00:31:26]
Step one --- there's a lot of steps and, as usual, I don't do all of them: \inlinemetanote{writes Step 1: What is the set $R/I$?} What is the set $R/I$? And here the notation is totally standard: everybody uses just this notation. What is the set?

If you remember linear algebra, that should help you now. We define an equivalence relation \inlinemetanote{writes We define an equivalence relation $\equiv$ on $R$} on $R$. Sometimes this equivalence relation could be denoted by this, \inlinemetanote{writes $\equiv_I$} but you have to carry around a lot of stuff, so I'm just going to leave that $I$ out. We're going to define an equivalence relation on $R$.

Here is the definition: if I have $x$ and $y$ in $R$, then $x$ is equivalent to $y$ \inlinemetanote{writes Def: $x, y \in R$, $x \equiv y$} --- and as I said, if you like, you can put $I$ here, because the equivalence relation depends on $I$ --- if and only if $x - y$ is in the ideal $I$. \inlinemetanote{writes $\iff x - y \in I$} That's the definition.

In words --- as an aside: \inlinemetanote{writes \qt{Side:}} if you want to be really precise, you could write $\equiv_I$, \inlinemetanote{writes $\equiv_I$} because this depends on $I$ --- this is usually put: \qt{$x$ is equivalent to $y \bmod I$}. \inlinemetanote{writes $x$ is equivalent to $y \bmod I$} This is very standard language: $x$ is equivalent to $y \bmod I$. That's what this equivalence relation is; it's just like for the integers.
\end{speech}

\begin{content}[The Equivalence Relation Modulo an Ideal]
\begin{definition}[Congruence Modulo an Ideal]\label[definition]{def:congruence-modulo-ideal}
Let $R$ be a ring and $I \trianglelefteq R$ an ideal. For elements $x, y \in R$, we define the relation of \newterm{congruence modulo $I$} (denoted $x \equiv_I y$ or simply $x \equiv y$) by:
\begin{equation}\label{eq:congruence-def}
x \equiv y \iff x - y \in I.
\end{equation}
In words, we say that \qt{$x$ is equivalent to $y$ modulo $I$}.
\end{definition}

\begin{steps}[Connection to Modular Arithmetic in $\mathbb{Z}$]
This directly generalizes congruence modulo $n$ in the integers:
\[
x \equiv y \pmod n \iff x - y \in n\mathbb{Z},
\]
where $I = n\mathbb{Z} = (n)$ is the principal ideal generated by $n$.
\end{steps}
\end{content}

\begin{speech}[00:31:26 - 00:33:28]
Formally, you now have to check that it's an equivalence relation. To do that, you have to remember the three properties of an equivalence relation: $x \equiv x$, which is zero being in the ideal; the second property is if $x \equiv y$, then $y \equiv x$, which would just put a minus sign here, but we know that $I$ is a subgroup under addition, so it's closed under taking negatives; and the last one is transitivity. Maybe I'll do the last one: the transitive property. \inlinemetanote{writes Transitivity:}

The transitive property says: if $x \equiv y$ and $y \equiv z$, then $x \equiv z$. \inlinemetanote{writes if $x \equiv y$ and $y \equiv z$ then $x \equiv z$} Now, $x \equiv y$ means $x - y \in I$, \inlinemetanote{writes $x - y \in I$} and $y \equiv z$ means $y - z \in I$. \inlinemetanote{writes $y - z \in I$} What I need to show is $x - z \in I$. But these two elements are in $I$, so I can add them: \inlinemetanote{writes $(x-y) + (y-z) \in I$} since $I$ is closed under addition, I can add them. Is that clear? And then, happily, the $y$ cancels out, so you get $x - z \in I$, \inlinemetanote{writes $x - z \in I$} and we're happy about that. That's the proof of transitivity.

If you think about what we've used: we've used zero is in $I$, we've used the minus sign is in $I$, and we've used that $I$ is closed under addition; so, we've used all the group conditions for $I$. We couldn't get away with less. What we have not used is the multiplicative condition --- that's going to come later. An ideal has two assumptions: one is the additive structure, it's a subgroup, and we've used it all here; but we have not touched the multiplicative condition. That'll come in a moment.
\end{speech}

\begin{content}[Verification: Equivalence Relation Properties]
\begin{proposition}[$\equiv_I$ is an Equivalence Relation]\label[proposition]{prop:equiv-rel-modulo-ideal}
Let $I \trianglelefteq R$ be an ideal. The relation $\equiv_I$ is an equivalence relation on $R$.
\end{proposition}

\begin{proof}
We verify the three defining axioms:
\begin{enumerate}[label=\textbf{(\arabic*)}]
    \item \textbf{Reflexivity ($x \equiv x$):}
    For any $x \in R$, $x - x = 0 \in I$ because $(I, +, 0)$ is a subgroup. Thus, $x \equiv x$.
    \item \textbf{Symmetry ($x \equiv y \implies y \equiv x$):}
    If $x \equiv y$, then $x - y \in I$. Since $(I, +, 0)$ contains additive inverses:
    \[
    y - x = -(x - y) \in I \implies y \equiv x.
    \]
    \item \textbf{Transitivity ($x \equiv y \land y \equiv z \implies x \equiv z$):}
    Suppose $x \equiv y$ and $y \equiv z$. By definition:
    \[
    x - y \in I \quad \text{and} \quad y - z \in I.
    \]
    Since $I$ is closed under addition, their sum lies in $I$:
    \[
    (x - y) + (y - z) \in I \implies x - z \in I \implies x \equiv z.
    \]
\end{enumerate}
\end{proof}

\begin{steps}[Subgroup Property Sufficiency]
Notice that verifying that $\equiv_I$ is an equivalence relation relies \textbf{exclusively} on the fact that $(I, +, 0)$ is an additive subgroup of $(R, +, 0)$. The multiplicative absorption condition $R \cdot I \subseteq I$ has not yet been utilized!
\end{steps}
\end{content}

\begin{hidden}
% timestamp: 00:33:28
% topic: Definition of congruence modulo an ideal and proof that it is an equivalence relation.
% board_state: def:congruence-modulo-ideal, prop:equiv-rel-modulo-ideal
% next_goal: Define the quotient set R/I as equivalence classes and examine naming conventions.
% open_loops: none
\end{hidden}

\begin{speech}[00:33:28 - 00:34:48]
If it's an equivalence relation, then you can imagine that we're going to take its equivalence classes. \inlinemetanote{erases the right portion of the lower board}

Now we can define the set: $R/I$ as a set is equal to the set of equivalence classes \inlinemetanote{writes $R/I = \text{set of equivalence classes}$} of this relation on $R$. \inlinemetanote{writes of $\equiv_I$ on $R$} So, I've defined this as a set.
\end{speech}

\begin{content}[The Quotient Set $R/I$]
\begin{definition}[The Quotient Set $R/I$]\label[definition]{def:quotient-set}
The underlying set of the quotient ring $R/I$ is defined as the set of all equivalence classes of $R$ with respect to the equivalence relation $\equiv_I$:
\begin{equation}\label{eq:quotient-set-def}
R/I := R / {\equiv_I} = \{ [x] \mid x \in R \}.
\end{equation}
\end{definition}
\end{content}

\subsection{Naming Conventions for Equivalence Classes (Cosets)}

\begin{speech}[00:34:48 - 00:36:05]
And now, the question of names: how do I name things? That's the same discussion we had about the rational numbers, the same for the field of fractions. One of the conceptual issues about the quotient is that elements do not have unique names. And this is just something you have to live with.

The question is: how do I name an element in $R/I$? \inlinemetanote{writes How do I name an element in $R/I$?}

One way to do it is: you take $x \in R$. \inlinemetanote{writes $x \in R$} Since $x$ is in $R$, it determines an equivalence class: the equivalence class which $x$ lies in. And that is usually called $\bar{x}$. \inlinemetanote{writes $\bar{x} \in R/I$ the equivalence class that $x$ lies in} This is the equivalence class that $x$ lies in. You take $x$, you get an equivalence class, and that's usually denoted $\bar{x}$. That's what we do with integers also. Is that okay?
\end{speech}

\begin{speech}[00:36:05 - 00:37:33]
There are three standard ways to denote this. If $x \in R$, we can write $\bar{x} \in R/I$. \inlinemetanote{writes $x \in R \implies \bar{x} \in R/I$}

If you really want to write things out explicitly, you can write it like that: $x + I \in R/I$, \inlinemetanote{writes $x + I \in R/I$} because the equivalence class of $x$ is equal to all those elements $x + i$ where $i \in I$. \inlinemetanote{writes $[x] = \{x + i \mid i \in I\}$} Because if you're equivalent to $x$, the difference is in $I$. That is explicitly what the equivalence class is. And since it's $x$ plus something in $I$, some people will write it like this: $x + I$. \inlinemetanote{points to $x + I$} It's kind of messy to keep that around, but in some arguments it's nice. That's the second way to write it.

The third way to write it --- which is by far the most common --- is: you just call it $x$ or $[x]$. \inlinemetanote{writes $[x]$ or just $x$} That's by far the most common. But if you use this, you have to be very careful: when $x$ was born, it was an element, and then when you put it here, it becomes an equivalence class; you have to keep track of the \qt{data type} of $x$ in your head. Those are the three ways to do it. They all have their own advantages --- well, this one has limited advantage, but okay.
\end{speech}

\begin{content}[Three Standard Notations for Residue Classes (Cosets)]
An element of $R/I$ is an equivalence class containing elements of $R$. Because representations are not unique (different representatives can specify the exact same class), three standard notations are commonly employed:

\begin{enumerate}[label=\textbf{(\arabic*)}]
    \item \textbf{Overline Notation ($\bar{x}$):}
    Denoted $\bar{x} \in R/I$. This is compact and mirrors the notation used for residue classes $\mathbb{Z}/n\mathbb{Z}$.
    \item \textbf{Coset Notation ($x + I$):}
    Explicitly writing the equivalence class as an affine shift of the ideal:
    \begin{equation}\label{eq:coset-explicit-def}
    [x] = x + I := \{ x + i \mid i \in I \} \subseteq R.
    \end{equation}
    Since $y \equiv x \iff y - x = i \in I \iff y = x + i$, the equivalence class of $x$ is literally the subset $x + I$.
    \item \textbf{Bare Representative Notation ($x$ or $[x]$):}
    Writing simply $[x]$ or, by abuse of notation, $x \in R/I$. While by far the most common in advanced algebra, this requires carefully tracking whether $x$ represents an individual ring element in $R$ or an entire equivalence class in $R/I$ (the \qt{data type} in one's head).
\end{enumerate}
\end{content}

\subsection{Step 2 \& 3: Distinguished Elements and the Well-Definedness of Addition}

\begin{speech}[00:37:33 - 00:39:08]
All right, so since we've defined the set, we need to now define addition. First, let's do zero and one: we define zero in $R/I$. \inlinemetanote{writes $0 \in R/I$, Def: $0 := \bar{0}$} By definition, $0 := \bar{0}$, because this is the equivalence class of zero. Zeros are always called zeros. So, that's the zero here: it's the equivalence class of zero, which is not a surprise. Similarly, the definition of one is that it's the equivalence class of one: $1 := \bar{1}$. \inlinemetanote{writes Def: $1 := \bar{1}$} Now we have at least this set $R/I$, and we have some distinguished elements: zero and one.

Now, let us define addition. \inlinemetanote{writes Let us define $+$} And it's still called plus. Very formally, you shouldn't call the binary operations in two different rings the same thing, because they're kind of different. But no one does that; everyone calls it plus, and it makes you have to concentrate to follow the definitions.

What does that mean? To define plus, I take an element of $R/I$ and I take another element of $R/I$, and I have to define this plus: \inlinemetanote{writes $\bar{x} + \bar{y} \overset{\text{def}}{=}$}
\end{speech}

\begin{content}[Distinguished Elements and the Definition of Addition]
\begin{definition}[Zero and Identity in $R/I$]\label[definition]{def:distinguished-elements}
The distinguished neutral elements in the quotient ring $R/I$ are defined by the classes of the respective identities in $R$:
\begin{align}
0_{R/I} &:= \bar{0} = 0 + I = I, \label{eq:zero-quotient-def} \\
1_{R/I} &:= \bar{1} = 1 + I. \label{eq:one-quotient-def}
\end{align}
\end{definition}

\begin{definition}[Addition on $R/I$]\label[definition]{def:addition-quotient}
For any two equivalence classes $\bar{x}, \bar{y} \in R/I$, their sum is defined by adding representatives in $R$:
\begin{equation}\label{eq:addition-quotient-formula}
\bar{x} + \bar{y} := \overline{x + y}, \quad \text{or equivalently} \quad (x + I) + (y + I) := (x + y) + I.
\end{equation}
\end{definition}
\end{content}

\begin{speech}[00:39:08 - 00:40:48]
This equation here is the definition. And how am I going to define it? Well, I only know how to add $x$ and $y$ one way, because they were born in $R$, and it's the only thing I know how to do. So, this is the only thing possibly I could write: $\overline{x + y}$. \inlinemetanote{writes $\overline{x + y}$} To read it, you have to concentrate so that you know what's being defined: the left-hand side is being defined, while the right-hand side was already previously known in $R$. This whole equation is the definition of plus.

And it's not a full definition yet, because you have to check that it's well-defined: these elements could have different names. Let's check that it's well-defined in $R/I$: \inlinemetanote{writes Let us check it is well-defined in $R/I$}

What could we do? We could have $\bar{x} = \bar{x}'$, where $x' = x + i$ with $i \in I$. \inlinemetanote{writes $\bar{x} = \bar{x}', x' = x + i, i \in I$} We could have just had a different name for the same equivalence class. And if we had a different name, we know it will just be some shift of $x$ by an element of $I$, because their difference is in $I$ by definition. Similarly, for $y$, we could have a different name: $y' = y + j$, where $j \in I$. \inlinemetanote{writes $\bar{y} = \bar{y}', y' = y + j, j \in I$}
\end{speech}

\begin{hidden}
% timestamp: 00:40:48
% topic: Defining addition in R/I and setting up the check for well-definedness.
% board_state: def:distinguished-elements, def:addition-quotient
% next_goal: Complete the proof of well-definedness of addition in R/I.
% open_loops: none
\end{hidden}

\begin{speech}[00:40:48 - 00:42:38]
Okay, so now we have to test that this definition is the same. We already have the formula for the names $x$ and $y$. Let's try to see what happens if we use these names: we'd better get the same answer! If we use these names, \inlinemetanote{writes $\bar{x}' + \bar{y}'$} this definition would give us $\overline{x' + y'}$. \inlinemetanote{writes $= \overline{x' + y'}$} Then we use the formulas for $x'$ and $y'$, and then I regroup them --- well, okay, maybe I don't: \inlinemetanote{writes $= \overline{(x+i) + (y+j)}$} $(x + i) + (y + j)$; and then I regroup it: $(x + y) + (i + j)$. \inlinemetanote{writes $= \overline{(x+y) + (i+j)}$} But now $i + j$ has to be in the ideal $I$, because $I$ is closed under addition. And so, this equivalence class is actually the same as that equivalence class, because it's the equivalence class of the same thing shifted by an element of $I$.

So, the answer is: we have checked that if I use different names and combine them with addition, we get the same equivalence class on the right-hand side. This proves it's well-defined. And if you remember, as I said, from linear algebra, you've done all of this before.

Then, formally, now that I've defined addition and zero, you have to formally check that this addition is associative, that zero plays the role of a zero, and that there are additive inverses; and I leave all of this to you. \inlinemetanote{erases the middle board}

But there is multiplication, and that's worth doing carefully, because so far we have not used the multiplicative part of the hypothesis on the ideal. And it comes out now, and it's kind of crucial. So, I will try to do this carefully, maybe after the break, but let me just say for...
\end{speech}

\begin{content}[Proof that Addition is Well-Defined on $R/I$]
\begin{proposition}[Well-Definedness of Addition]\label[proposition]{prop:addition-well-defined}
The addition operation defined in \cref{def:addition-quotient} is well-defined: it is independent of the choice of representatives.
\end{proposition}

\begin{proof}
Suppose $\bar{x} = \bar{x}'$ and $\bar{y} = \bar{y}'$ in $R/I$. By \cref{def:congruence-modulo-ideal}, there exist elements $i, j \in I$ such that:
\begin{align}
x' &= x + i, \label{eq:x-prime-shift} \\
y' &= y + j. \label{eq:y-prime-shift}
\end{align}
Evaluating the sum using the primed representatives yields:
\begin{align*}
\bar{x}' + \bar{y}' &= \overline{x' + y'} \\
&= \overline{(x + i) + (y + j)} \\
&= \overline{(x + y) + (i + j)}.
\end{align*}
Since $(I, +, 0)$ is an additive subgroup, $i, j \in I$ implies $i + j \in I$. Therefore, $(x + y) + (i + j) \equiv (x + y) \pmod I$, which implies:
\[
\overline{(x + y) + (i + j)} = \overline{x + y} = \bar{x} + \bar{y}.
\]
Thus, the sum is independent of representatives.
\end{proof}

\begin{steps}[The Quotient Additive Group]
Since addition is well-defined, associativity, commutativity, the identity $\bar{0}$, and inverses $-\bar{x} = \overline{-x}$ all inherit directly from $(R, +, 0)$. Hence, $(R/I, +, \bar{0})$ forms an abelian group.
\end{steps}
\end{content}

\subsection{Step 4: Multiplication and the Crucial Role of the Ideal Absorption Property}

\begin{speech}[00:42:38 - 00:43:48]
The next step is to define multiplication on $R/I$. \inlinemetanote{writes Define $\cdot$ on $R/I$} That means, if we have $x$ and $y$, it's exactly as before: we need to define this, \inlinemetanote{writes $\bar{x} \cdot \bar{y} \overset{\text{def}}{=}$} and this is going to be the definition. And how are you going to define it? You only know how to do one thing to define things, so you really have no choice except this: \inlinemetanote{writes $\overline{xy}$} $\overline{x y}$. I could put the dot there too if I want.

This just says that if I want to multiply two equivalence classes, I take one element of one equivalence class, another element of the other equivalence class, I multiply them in the ring, and then take the resulting equivalence class. Is that clear?
\end{speech}

\begin{content}[Definition of Multiplication on $R/I$]
\begin{definition}[Multiplication on $R/I$]\label[definition]{def:multiplication-quotient}
For equivalence classes $\bar{x}, \bar{y} \in R/I$, their product is defined by multiplying representatives in $R$:
\begin{equation}\label{eq:multiplication-quotient-def}
\bar{x} \cdot \bar{y} := \overline{x \cdot y}, \quad \text{or equivalently} \quad (x + I) \cdot (y + I) := (x \cdot y) + I.
\end{equation}
\end{definition}
\end{content}

\begin{speech}[00:43:48 - 00:45:11]
And then I have to prove this is well-defined. \inlinemetanote{writes Prove well-defined} So, I use the same choices: suppose I pick these different names and I use the definition, \inlinemetanote{writes $\bar{x}' \cdot \bar{y}' = \overline{x' y'}$} then I get $\overline{x' y'}$, and then I use their formulas: \inlinemetanote{writes $= \overline{(x+i)(y+j)}$} $(x + i)(y + j)$ \inlinemetanote{corrects to $(x+i)(y+j)$}.

And when I multiply them out --- this is the crucial part ---, I get... sorry, there's no primes here. I get $x y$, and then I get $i \cdot y + j \cdot x$ --- it's commutative, so I can write it either way --- plus $i j$. \inlinemetanote{writes $= \overline{xy + iy + jx + ij}$ feedback on terms}

With one name, I get $\overline{x y}$. With the other name, I get this. And the difference is this. So, the only way this is well-defined is if this whole thing is in $I$. \inlinemetanote{brackets $iy + jx + ij$ and writes $\in I$} That's the only way it could be well-defined. And it is in $I$! Here, finally, we use the fact that if you have an element of $I$ and multiply it by anything, you stay in $I$. So, that's why it's here. Okay, so...
\end{speech}

\begin{content}[Proof that Multiplication is Well-Defined on $R/I$]
\begin{proposition}[Well-Definedness of Multiplication]\label[proposition]{prop:multiplication-well-defined}
The multiplication operation defined in \cref{def:multiplication-quotient} is well-defined.
\end{proposition}

\begin{proof}
Let $\bar{x}' = \bar{x}$ and $\bar{y}' = \bar{y}$. Then $x' = x + i$ and $y' = y + j$ for some $i, j \in I$. Expanding the product of the alternative representatives:
\begin{align*}
\bar{x}' \cdot \bar{y}' &= \overline{x' \cdot y'} \\
&= \overline{(x + i)(y + j)} \\
&= \overline{x y + i y + j x + i j}.
\end{align*}
To establish that $\overline{(x + i)(y + j)} = \overline{x y}$, we must verify that the difference belongs to $I$:
\[
(x y + i y + j x + i j) - x y = i y + j x + i j \stackrel{?}{\in} I.
\]
We examine each of the three cross-terms:
\begin{itemize}
    \item $i \in I$ and $y \in R \implies i \cdot y \in I$ by the \textbf{multiplicative absorption property} of ideals ($I \cdot R \subseteq I$).
    \item $j \in I$ and $x \in R \implies j \cdot x \in I$ by the \textbf{multiplicative absorption property} ($R \cdot I \subseteq I$).
    \item $i, j \in I \implies i \cdot j \in I$ (in fact, $I \cdot I \subseteq I$).
\end{itemize}
Since $(I, +, 0)$ is closed under addition:
\[
i y + j x + i j \in I \implies \overline{x y + i y + j x + i j} = \overline{x y} = \bar{x} \cdot \bar{y}.
\]
Therefore, multiplication on $R/I$ is strictly well-defined.
\end{proof}

\begin{steps}[Why Subrings Fail and Ideals are Necessary]
If $I$ were merely an additive subgroup or a subring (closed under multiplication within itself, but not absorbing arbitrary elements of $R$), the terms $i y$ and $j x$ would generally fail to lie in $I$, causing the multiplication on equivalence classes to collapse! \textbf{This step is the sole reason why ideals, rather than general subrings, are required to define quotient rings.}
\end{steps}
\end{content}

\begin{hidden}
% timestamp: 00:45:11
% topic: Well-definedness of multiplication in R/I and mid-lecture break.
% board_state: def:multiplication-quotient, prop:multiplication-well-defined
% next_goal: Recap of ideal axioms and why absorption is crucial after the break.
% open_loops: none
\end{hidden}

\begin{pause}[Mid-Lecture Break]
The lecturer takes a short pause for the mid-lecture break. The lecture resumes at timestamp 00:45:11 with a pedagogical recap of the essential role played by the ideal axioms.
\end{pause}

\subsection{Recap: The Crucial Role of the Ideal Axioms}

\begin{speech}[00:45:11 - 00:46:13]
Okay, so we start, and the last moment in the first half of the class was about this observation, and it's a really important thing. So, I repeat it now: this is the only time...

If you remember, the definition of an ideal had two clauses. The first is that $(I, +, 0)$ is an additive subgroup, and we've used that repeatedly: in the definition of the equivalence relation, in the definition of addition; and if you wanted to prove associativity and all of those properties that I didn't prove, you'd be using it here.

But it has a second property, which was the multiplicative property: that for every element of the ring and for every $a$ in the ideal, the product is in the ideal. \inlinemetanote{writes on the lower right board: Def of Ideal: (1) $(I, +, 0) < (R, +, 0)$ subgroup, (2) $\forall r \in R, \forall a \in I \implies ra \in I$} That's just a totally separate property, because the first one doesn't have anything to do with multiplication. This one has to do with multiplication, and this is the first time we use it.
\end{speech}

\begin{content}[Review: The Two Defining Axioms of an Ideal]
\begin{definition}[Ideal Axioms Reviewed]\label[definition]{def:ideal-axioms-reviewed}
A subset $I \subseteq R$ is an ideal if and only if:
\begin{enumerate}[label=\textbf{(\arabic*)}]
    \item \textbf{Additive Subgroup Axiom:}
    \[
    (I, +, 0) \le (R, +, 0) \quad \text{is an additive subgroup.}
    \]
    \item \textbf{Multiplicative Absorption Axiom:}
    \[
    \forall r \in R, \quad \forall a \in I \implies r \cdot a \in I.
    \]
\end{enumerate}
\end{definition}

\begin{steps}[Division of Labor Between the Two Axioms]
\begin{itemize}
    \item Axiom \textbf{(1)} governs all additive aspects: transitivity of equivalence classes, well-definedness of addition, additive identity, and additive inverses in $R/I$.
    \item Axiom \textbf{(2)} is reserved solely for multiplication: ensuring that products of cosets do not depend on the representative chosen.
\end{itemize}
\end{steps}
\end{content}

\begin{speech}[00:46:13 - 00:47:25]
And it's used to show that this definition of multiplication --- which, as I tried to say, philosophically you didn't have very much choice about: if you want to multiply two equivalence classes, you pick an element in one equivalence class, you pick an element in the other, you multiply those elements, and you take the equivalence class; that's pretty much the only choice you have to define it, because you start with just an abstract ring.

And to show this is well-defined, which I did here quickly at the end with different elements in those equivalence classes, you come to this issue: you're taking the equivalence class of this element, and it has to be the same as the equivalence class of that element. And the only way it is the same is if the extra terms, when you combine them, are all in $I$.

And to see that: this is a product of two elements of $I$. That's okay, because for any element of $R$ and any element of $I$, you only need one element of $I$. So, this one is in $I$; and these two have one element of $I$ and one element of $R$ each, because $x$ and $y$ are arbitrary elements of $R$. So, these are three times we apply this rule.
\end{speech}

\subsection{Completion of the Quotient Ring Structure}

\begin{speech}[00:47:25 - 00:48:47]
Okay, so that proves...

And this is one aspect of the construction which is not in your linear algebra class. Although it's not far from what happened in your linear algebra class when you considered the problem of scalar multiplication on the quotient, which I won't repeat, but anyway. So, even that shouldn't be totally foreign to you.

So, now we've constructed this: \inlinemetanote{writes $((R/I), +, \cdot, \bar{0}, \bar{1})$ on the lower left board} we have plus, we have multiplication, we have $0$ and $1$. \inlinemetanote{adds checkmarks above the symbols}

The tricky part --- it's not even that tricky, but the part that requires the foundational work --- is the definition as equivalence classes, and making sure these are well-defined. For the rest, you just follow your nose. The identity: you can imagine, you have to check $\bar{1} \cdot \bar{x}$, \inlinemetanote{writes $\bar{1} \cdot \bar{x} = \overline{1 \cdot x} = \bar{x}$ with a checkmark} so by definition this is $\overline{1 \cdot x}$, and so that's $\bar{x}$. That checks one of the identity properties, and there are about $10$ other properties you have to check, and they're all checked like this.

So, that's the construction, and that's the main construction this week. Last week we had the construction of the field of fractions, which you should understand completely; and now there is this construction, which is the construction of the quotient ring, which is your task to understand really completely this week: the construction.
\end{speech}

\begin{content}[The Complete Quotient Ring Structure]
\begin{theorem}[The Quotient Ring $(R/I, +, \cdot, \bar{0}, \bar{1})$]\label[theorem]{thm:quotient-ring-structure}
Let $R$ be a commutative ring with identity, and let $I \trianglelefteq R$ be an ideal. The quintuple:
\begin{equation}\label{eq:quotient-ring-quintuple}
(R/I, +, \cdot, \bar{0}, \bar{1})
\end{equation}
equipped with the coset operations:
\[
\bar{x} + \bar{y} := \overline{x + y}, \qquad \bar{x} \cdot \bar{y} := \overline{x \cdot y}, \qquad 0_{R/I} := \bar{0}, \qquad 1_{R/I} := \bar{1},
\]
forms a commutative ring with identity.
\end{theorem}

\begin{proof}
Having proven that addition and multiplication are well-defined, the remaining ring axioms follow directly from the ring axioms of $R$ applied to representatives:
\begin{itemize}
    \item \textbf{Multiplicative Identity:} $\bar{1} \cdot \bar{x} = \overline{1 \cdot x} = \bar{x}$ for all $\bar{x} \in R/I$.
    \item \textbf{Associativity:} $(\bar{x} \cdot \bar{y}) \cdot \bar{z} = \overline{(x y) z} = \overline{x (y z)} = \bar{x} \cdot (\bar{y} \cdot \bar{z})$.
    \item \textbf{Distributivity:} $\bar{x} \cdot (\bar{y} + \bar{z}) = \overline{x (y + z)} = \overline{x y + x z} = (\bar{x} \cdot \bar{y}) + (\bar{x} \cdot \bar{z})$.
    \item \textbf{Commutativity:} $\bar{x} \cdot \bar{y} = \overline{x y} = \overline{y x} = \bar{y} \cdot \bar{x}$.
\end{itemize}
\end{proof}
\end{content}

\begin{hidden}
% timestamp: 00:48:47
% topic: Completion of the quotient ring construction R/I.
% board_state: def:multiplication-quotient, prop:multiplication-well-defined, thm:quotient-ring-structure
% next_goal: Address student question regarding relationship of lectures to textbook and exam scope.
% open_loops: none
\end{hidden}

\begin{speech}[00:48:47 - 00:51:00]
Are there any questions about this?

People were asking me: what's the relationship of this class to the book? That's a really good question.

Every week I pick out some pages in that book --- it's a huge book ---, and somehow the class covers a very small part of the book, so it's not your goal to read that book from cover to cover, although it's not a bad thing to do!

Every week I pick out some topics there, and my lectures are roughly on those topics. Now, my notation is sometimes different, and the presentation is sometimes different in the book. The book is useful if you want... I skip some details in lectures; if you want all the details, sometimes they're in the book. That selection for the whole semester also covers the actual curriculum I'm required to cover, more or less, for this course, so it's helpful for me. But I would say: use the book if it's helpful for you; if not, you can use the lectures.

What happened when I taught this course at ETH in 2018 was that some students wrote some notes, and they're really great notes. Of course, they sent them to me afterwards, but it's $8$ years ago and I couldn't find them. But probably some of your colleagues write some good notes; you can also use that. The lectures are recorded, so you can also look at that.

Another option is that I gave some really good references --- some really classic algebra references. They have the unfortunate quality that their PDF file is not available through the library, but you can look at those books also. But if you look at the other books, the material won't be exactly in the order I'm giving you. If you want a book where it's in the order I'm giving you, then you should follow those page numbers in the actual reference; and that book is freely available on ITS, on the ETH website. So, that's the discussion about the book: use it if it's helpful.

For the exam material and things, I will only use material that I lectured on in class, not some exotic topic in the book that I didn't lecture on.

Okay.
\end{speech}

\begin{note}[Course References and Exam Scope]
The lecturer answers student inquiries regarding the course literature:
\begin{itemize}
    \item \textbf{Textbook Alignment:} Weekly reading assignments correspond to specific sections of the comprehensive course textbook (freely accessible to students via the ETH website). The textbook provides supplementary proofs and alternative presentations.
    \item \textbf{Exam Scope:} The examinations will strictly test the concepts, definitions, and theorems presented during the lectures, rather than extraneous textbook topics.
\end{itemize}
\end{note}

\section{The First Isomorphism Theorem for Rings}

\begin{speech}[00:51:00 - 00:52:37]
All right, so we go back to this issue now. We've made some progress. \inlinemetanote{writes $f \colon R \to S$, $R, S$ rings, $f$ hom on the upper middle board} If I have $f \colon R \to S$ --- where $R$ and $S$ are rings and $f$ is a homomorphism ---, then this is sometimes called the First Isomorphism Theorem. \inlinemetanote{writes First Isom. Theorem}

And it plays the role of the missing piece that we had in linear algebra that I wanted to develop here: that says that this ring, $R/\ker(f)$ \inlinemetanote{writes $R/\ker(f)$} --- because the kernel is an ideal, and if we have an ideal we can take the quotient, which is what the last $45$ minutes were about ---, this ring is canonically isomorphic to the image: $\cong \operatorname{Im}(f)$, \inlinemetanote{writes $\cong \operatorname{Im}(f)$} okay?

That plays the role of the statement in vector spaces. Is that clear? That's the replacement. So, we have to prove this.

Let us see what we have to do: we start with $f \colon R \to S$, \inlinemetanote{writes We start $f \colon R \to S$} and we will define a homomorphism $\bar{f} \colon R/\ker(f) \to S$. \inlinemetanote{writes We will define a hom $\bar{f} \colon R/\ker(f) \to S$} Just to start. And how do you think we're going to define this?
\end{speech}

\begin{content}[Statement of the First Isomorphism Theorem for Rings]
\begin{theorem}[First Isomorphism Theorem for Rings]\label[theorem]{thm:first-isomorphism-theorem-statement}
Let $f \colon R \to S$ be a homomorphism of commutative rings. Then the kernel $\ker(f)$ is an ideal of $R$, the image $\operatorname{Im}(f)$ is a subring of $S$, and there exists a canonical ring isomorphism:
\begin{equation}\label{eq:first-isomorphism-ring}
\bar{f} \colon R/\ker(f) \overset{\cong}{\longrightarrow} \operatorname{Im}(f), \qquad \bar{x} \mapsto f(x).
\end{equation}
In particular, if $f$ is surjective, then $R/\ker(f) \cong S$.
\end{theorem}
\end{content}

\begin{note}[Board Cleaning Transition]
\inlinemetanote{The lecturer washes the lower blackboard completely with a damp sponge from 00:52:37 to 00:53:30 to prepare the space for the three-step proof of the First Isomorphism Theorem.}
\end{note}

\begin{proofWrapper}[Proof of the First Isomorphism Theorem]

\subsection{Step 1: Construction and Well-Definedness of the Induced Map \texorpdfstring{$\bar{f}$}{fbar}}

\begin{speech}[00:53:30 - 00:54:07]
Okay, so we have to define $\bar{f}$. What does that mean? We have to define this: \inlinemetanote{writes $\bar{f}(\bar{x}) \overset{\text{def}}{=} f(x) \in S$ on the lower board} $\bar{f}$ is a homomorphism that starts from $R/\ker(f)$. What this thing eats is elements of $R/\ker(f)$ \inlinemetanote{writes $\bar{x} \in R/\ker(f)$} --- it eats equivalence classes.

What do you think we're going to do here? Yeah?
\end{speech}

\begin{student}[Student Answer]
Send it to $f(x)$?
\end{student}

\begin{speech}[00:54:08 - 00:54:34]
Yeah. In this kind of generality, there's very little you can do, because you can imagine $R$ is like one country, $S$ is another, and there's only one flight between them, and you've got to take it! There's just no other way. To get in here, you've got to get on $f$ somehow, and that's the only way. So, that's the definition.
\end{speech}

\begin{speech}[00:54:34 - 00:55:47]
When you define it like this, as I said, it's the standard pain about this naming business. And, you know, I say it slightly facetiously, but naming is an incredibly crucial thing. As you study mathematics more and more, you'll see this idea of what you can name, what you can name precisely, and how you can name it turns out to be an incredibly important, subtle point. So, it's good to get used to thinking about that.

Anyway, the point is: we have to show this is well-defined. \inlinemetanote{writes Show $\bar{f}$ is well defined}

What does that mean? Suppose we have somebody else in the equivalence class: $x' = x + i$, where $i \in \ker(f)$. \inlinemetanote{writes $x' = x + i, i \in \ker(f)$} If we take another representative in the same equivalence class, $\bar{x}' = \bar{x}$, \inlinemetanote{writes $\bar{x}' = \bar{x}$} it satisfies this equation. Then we must check that $\bar{f}(\bar{x}') = \bar{f}(\bar{x})$. \inlinemetanote{writes We must check $\bar{f}(\bar{x}') = \bar{f}(\bar{x})$} We have to check it.
\end{speech}

\begin{speech}[00:55:47 - 00:56:42]
Let's see: what is this? Here we use the definition: that's $f(x')$, \inlinemetanote{writes $\bar{f}(\bar{x}') \overset{\text{def}}{=} f(x')$} and then we use the definition of $x'$, so that's $f(x + i)$. \inlinemetanote{writes $= f(x+i)$}

I should have said here very carefully that $i \in \ker(f)$. \inlinemetanote{writes $i \in \ker(f)$} What do we do here? Well, we can expand this as $f(x) + f(i)$, \inlinemetanote{writes $= f(x) + f(i)$} because $f$ itself is a homomorphism. And now $i$ is in the kernel, so this is equal to $f(x)$, \inlinemetanote{writes $= f(x)$} because $f(i) = 0$, \inlinemetanote{writes $f(i) = 0$} and that's what we wanted, right? And this is equal to $\bar{f}(\bar{x})$. \inlinemetanote{writes $= \bar{f}(\bar{x})$} Okay.
\end{speech}

\begin{student}[Student Question]
Don't you need to first show that $\bar{f}$ is a homomorphism, because we are using the argument that $\bar{f}$ of $i$ bar, which is $\bar{f}$ of $0$ bar is $0$, but...
\end{student}

\begin{speech}[00:56:55 - 00:57:33]
No, I'm just showing it's well-defined: I only want to show this is equal to that.

All this $f$ stuff is fine, because $f$ is a homomorphism. Remember: $f$ is a homomorphism.

We have to do all the things you're saying, but there's a certain order you do it in. First, I define $\bar{f}$ --- this is part of the definition of $\bar{f}$. I propose a definition here. The problem with this definition is that it may depend on which element of the equivalence class I pick. And now I check that it doesn't depend on that. But I haven't proven this is a homomorphism yet; I will have to do that.
\end{speech}

\begin{student}[Student Question]
Where do we know that $i$ is in the kernel?
\end{student}

\begin{speech}[00:57:37 - 00:57:48]
Well, that's right here: for another representative, that's what it means, that their difference is in the kernel.
\end{speech}

\begin{speech}[00:57:48 - 00:58:45]
If that was confusing, I'll say it again, because it's kind of important.

I proposed this as a definition, but it has the flaw that I have to check that this definition does not depend on which element of the equivalence class I chose. And this is another element of the equivalence class. Any two elements of the equivalence class differ by an element of the kernel. So, that's the equation there, and I use this equation to show this map is well-defined.

Actually, I wrote out all the steps. In my experience in teaching, it doesn't always help writing out all the steps: there's a lot of steps, and every time you have to check what I'm using at every equality. You have to do that: I already know, but you have to verify it, and if I just write it, it doesn't always help.

So, that was a little painful. But we're not done with the pain!
\end{speech}

\begin{content}[Definition of $\bar{f}$ and Proof of Well-Definedness]
We define the induced map:
\begin{equation}\label{eq:induced-map-def}
\bar{f} \colon R/\ker(f) \longrightarrow S, \qquad \bar{f}(\bar{x}) := f(x).
\end{equation}

\begin{claim}[$\bar{f}$ is Well-Defined]\label[claim]{clm:fbar-well-defined}
The map $\bar{f}$ is independent of the choice of representative in the coset $\bar{x}$.
\end{claim}

\begin{proof}
Suppose $\bar{x}' = \bar{x}$ in $R/\ker(f)$. By definition of congruence modulo $\ker(f)$, there exists an element $i \in \ker(f)$ such that:
\[
x' = x + i.
\]
Applying the definition of $\bar{f}$ and the additive homomorphism property of $f$:
\begin{align*}
\bar{f}(\bar{x}') &= f(x') && \text{(by definition of $\bar{f}$)} \\
&= f(x + i) && \text{(since $x' = x + i$)} \\
&= f(x) + f(i) && \text{(since $f$ is a ring homomorphism)} \\
&= f(x) + 0_S && \text{(since $i \in \ker(f) \implies f(i) = 0_S$)} \\
&= f(x) \\
&= \bar{f}(\bar{x}).
\end{align*}
Thus, $\bar{f}(\bar{x}') = \bar{f}(\bar{x})$, proving that $\bar{f}$ is well-defined as a set-theoretic function.
\end{proof}
\end{content}

\begin{speech}[00:58:45 - 01:00:22]
This just shows it's well-defined, so we have $\bar{f} \colon R/\ker(f) \to S$ as a function now. \inlinemetanote{writes $\bar{f} \colon R/\ker(f) \to S$ function}

And we have to show all the properties to show it's a ring homomorphism.

We would have to show, for example, $\bar{f}(\bar{0})$: \inlinemetanote{writes $\bar{f}(\bar{0})$} what is this? We use our definition, and that's equal to $f(0)$, \inlinemetanote{writes $= f(0)$} and $f(0) = 0$ \inlinemetanote{writes $= 0$} because $f$ itself is a homomorphism. So, that's good: it's sending $0$ to $0$. That's maybe what you were worried about; now we've shown it.

We also have to show $\bar{f}(\bar{x} + \bar{y}) = \bar{f}(\bar{x}) + \bar{f}(\bar{y})$. \inlinemetanote{writes $\bar{f}(\bar{x}+\bar{y}) = \bar{f}(\bar{x}) + \bar{f}(\bar{y})$} You just have to expand these sides completely, and if you do that, you'll get a chain more or less similar to that, and you'll find that it's compatible with the definition.

So, expand it yourself! \inlinemetanote{writes expand yourself} As I said, you'll get a chain...
\end{speech}

\begin{speech}[01:00:22 - 01:02:03]
That's the additive structure. Then you also have to show that $\bar{f}$ respects the multiplicative structure: you have to check that $\bar{f}(\bar{x} \cdot \bar{y}) = \bar{f}(\bar{x}) \cdot \bar{f}(\bar{y})$. \inlinemetanote{writes check $\bar{f}(\bar{x} \cdot \bar{y}) = \bar{f}(\bar{x}) \cdot \bar{f}(\bar{y})$ on the lower left board}

If we do that: first, we have the definition of the product in the quotient, that's $\bar{f}(\overline{x y})$; \inlinemetanote{writes $\bar{f}(\overline{xy})$} and then, by definition of $\bar{f}$, that's $f(x y)$; \inlinemetanote{writes $= f(xy)$} and expanding this, that's $f(x) f(y)$. \inlinemetanote{writes $= f(x) f(y)$} And that's true because $f$ is a homomorphism --- true because $f$ was born a homomorphism! \inlinemetanote{writes True because $f$ is hom $R \to S$} Okay.

You have to concentrate with the bars, but really, the proof of this is just unraveling the definitions.
\end{speech}

\begin{content}[Proof that $\bar{f}$ is a Ring Homomorphism]
\begin{claim}[$\bar{f}$ is a Ring Homomorphism]\label[claim]{clm:fbar-is-homomorphism}
The well-defined map $\bar{f} \colon R/\ker(f) \to S$ preserves zero, identity, addition, and multiplication.
\end{claim}

\begin{proof}
Using the coset definitions in $R/\ker(f)$ and the ring homomorphism properties of $f \colon R \to S$:
\begin{enumerate}[label=\textbf{(\arabic*)}]
    \item \textbf{Preservation of Zero:}
    \[
    \bar{f}(\bar{0}) = f(0_R) = 0_S.
    \]
    \item \textbf{Preservation of Identity:}
    \[
    \bar{f}(\bar{1}) = f(1_R) = 1_S.
    \]
    \item \textbf{Preservation of Addition:}
    \begin{align*}
    \bar{f}(\bar{x} + \bar{y}) &= \bar{f}(\overline{x + y}) && \text{(by definition of addition in $R/\ker(f)$)} \\
    &= f(x + y) && \text{(by definition of $\bar{f}$)} \\
    &= f(x) + f(y) && \text{(since $f$ is a ring homomorphism)} \\
    &= \bar{f}(\bar{x}) + \bar{f}(\bar{y}). && \text{(by definition of $\bar{f}$)}
    \end{align*}
    \item \textbf{Preservation of Multiplication:}
    \begin{align*}
    \bar{f}(\bar{x} \cdot \bar{y}) &= \bar{f}(\overline{x \cdot y}) && \text{(by definition of multiplication in $R/\ker(f)$)} \\
    &= f(x \cdot y) && \text{(by definition of $\bar{f}$)} \\
    &= f(x) \cdot f(y) && \text{(since $f$ is a ring homomorphism)} \\
    &= \bar{f}(\bar{x}) \cdot \bar{f}(\bar{y}). && \text{(by definition of $\bar{f}$)}
    \end{align*}
\end{enumerate}
Therefore, $\bar{f}$ is a homomorphism of commutative rings.
\end{proof}
\end{content}

\begin{hidden}
% timestamp: 01:02:03
% topic: Well-definedness of induced map \bar{f} and verification that it is a ring homomorphism (Step 1).
% board_state: thm:first-isomorphism-theorem-statement, clm:fbar-well-defined, clm:fbar-is-homomorphism
% next_goal: Step 2: Surjectivity onto the image of f; Step 3: Injectivity via kernel of \bar{f}.
% open_loops: none
\end{hidden}

\subsection{Step 2: Surjectivity of \texorpdfstring{$\bar{f}$}{fbar} onto \texorpdfstring{$\operatorname{Im}(f)$}{Im(f)}}

\begin{speech}[01:02:03 - 01:03:46]
That's the first step in the proof. So, I remind you: we have $f \colon R \to S$, the ideal is $\ker(f)$, and we've defined $\bar{f} \colon R/\ker(f) \to S$. \inlinemetanote{writes $f \colon R \to S$, $\ker(f) \trianglelefteq R$, $\bar{f} \colon R/\ker(f) \to S$ is hom} And this is a ring homomorphism. We know the domain is a ring --- that was the first part of this class ---, we've defined $\bar{f}$, and we proved it's a ring homomorphism.

All right, step two: what is the image of $\bar{f}$? \inlinemetanote{writes Step 2: What is $\operatorname{Im}(\bar{f})$?} Where do we go here? The answer is: it's the same as the image of $f$. So, actually, $\bar{f}$ goes from $R/\ker(f)$ to the image of $f$. \inlinemetanote{writes $\bar{f} \colon R/\ker(f) \to \operatorname{Im}(f)$} This is some subring of $S$.

Why? Because the definition just says that if I take an equivalence class and choose a representative, I go to $f(x)$. So, everything I get is in the image of $f$. On the other hand, if I'm in the image of $f$, that means that I come from some element of $R$, and I can take its equivalence class. So, I get exactly the image of $f$.

That's step two, and this is a surjection. So, $\bar{f}$ is surjective onto the image of $f$, \inlinemetanote{writes $\bar{f}$ is surjective onto the image of $f$} okay?
\end{speech}

\begin{speech}[01:03:46 - 01:04:46]
Sometimes in math you write two arrows for surjectivity; \inlinemetanote{draws a two-headed arrow on the board} do you follow that kind of stuff or no? Okay, I won't give you more notation now.

First, we defined $\bar{f}$, and $\bar{f}$ is defined in the only possible way you can. The second step is: we realize it's surjective onto the image of $f$.
\end{speech}

\begin{content}[Surjectivity of the Induced Homomorphism onto the Image]
\begin{claim}[Surjectivity of $\bar{f}$ onto $\operatorname{Im}(f)$]\label[claim]{clm:fbar-surjective}
The image of $\bar{f}$ is precisely the subring $\operatorname{Im}(f) \subseteq S$. Consequently, the restricted codomain map:
\begin{equation}\label{eq:fbar-surjective-map}
\bar{f} \colon R/\ker(f) \twoheadrightarrow \operatorname{Im}(f)
\end{equation}
is a surjective ring homomorphism.
\end{claim}

\begin{proof}
We verify mutual inclusion of sets:
\begin{itemize}
    \item \textbf{$\operatorname{Im}(\bar{f}) \subseteq \operatorname{Im}(f)$:} For any coset $\bar{x} \in R/\ker(f)$, $\bar{f}(\bar{x}) = f(x) \in \operatorname{Im}(f)$ by definition of the image of $f$.
    \item \textbf{$\operatorname{Im}(f) \subseteq \operatorname{Im}(\bar{f})$:} Let $y \in \operatorname{Im}(f)$. By definition of the image, there exists an element $x \in R$ such that $f(x) = y$. Considering the coset $\bar{x} \in R/\ker(f)$, we have:
    \[
    \bar{f}(\bar{x}) = f(x) = y.
    \]
    Thus, $y \in \operatorname{Im}(\bar{f})$.
\end{itemize}
Hence, $\operatorname{Im}(\bar{f}) = \operatorname{Im}(f)$, proving that $\bar{f}$ maps $R/\ker(f)$ surjectively onto the subring $\operatorname{Im}(f)$.
\end{proof}
\end{content}

\subsection{Step 3: Injectivity and Trivial Kernel of \texorpdfstring{$\bar{f}$}{fbar}}

\begin{speech}[01:04:46 - 01:06:32]
And the third step --- this one requires some real concentration --- is: what is the kernel of $\bar{f}$? \inlinemetanote{writes Step 3: What is the $\ker(\bar{f})$?}

Let's try to think about that. Suppose I have $\bar{x} \in R/\ker(f)$, \inlinemetanote{writes Suppose $\bar{x} \in R/\ker(f)$} and $\bar{f}(\bar{x}) = 0 \in S$. \inlinemetanote{writes and $\bar{f}(\bar{x}) = 0 \in S$} This is another way of saying $\bar{x} \in \ker(\bar{f})$. \inlinemetanote{writes $(\bar{x} \in \ker(\bar{f}))$}

What can I conclude? Well, by definition, $\bar{f}(\bar{x}) = f(x)$, so this implies that $f(x) = 0 \in S$, \inlinemetanote{writes Then $f(x) = 0 \in S$} which implies $x \in \ker(f)$. \inlinemetanote{writes $\implies x \in \ker(f)$}

What have we done here? We said: we start with any equivalence class; if it's in the kernel of $\bar{f}$, then that element has to be in the kernel of $f$. And that means that equivalence class is the zero equivalence class, because it's equivalent to zero modulo the kernel: $x \equiv_{\ker(f)} 0$. \inlinemetanote{writes $\implies x \equiv_{\ker(f)} 0$} And this means that $\bar{x} = \bar{0} \in R/\ker(f)$. \inlinemetanote{writes $\implies \bar{x} = \bar{0} \in R/\ker(f)$}

So, that proves that the kernel of $\bar{f}$ is zero! That's step three: the conclusion is that the kernel of $\bar{f}$ is zero.
\end{speech}

\begin{content}[Injectivity of $\bar{f}$ via Triviality of the Kernel]
\begin{claim}[$\bar{f}$ is Injective]\label[claim]{clm:fbar-injective}
The kernel of the homomorphism $\bar{f} \colon R/\ker(f) \to \operatorname{Im}(f)$ is the zero ideal of the quotient ring:
\begin{equation}\label{eq:kernel-fbar-zero}
\ker(\bar{f}) = (\bar{0}).
\end{equation}
Consequently, $\bar{f}$ is injective.
\end{claim}

\begin{proof}
Let $\bar{x} \in \ker(\bar{f}) \subseteq R/\ker(f)$. By definition of the kernel:
\[
\bar{f}(\bar{x}) = 0_S.
\]
Using the definition of the map $\bar{f}(\bar{x}) := f(x)$, this means:
\[
f(x) = 0_S \implies x \in \ker(f).
\]
By definition of cosets modulo $\ker(f)$, an element belongs to the ideal if and only if its coset coincides with the zero coset:
\[
x \in \ker(f) \iff x \equiv_{\ker(f)} 0 \iff \bar{x} = \bar{0} \in R/\ker(f).
\]
Hence, the only element in $\ker(\bar{f})$ is the zero coset $\bar{0}$, establishing:
\[
\ker(\bar{f}) = (\bar{0}).
\]
By the injectivity criterion established earlier in \cref{prop:injective-kernel}, a ring homomorphism is injective if and only if its kernel is trivial. Therefore, $\bar{f}$ is injective.
\end{proof}

\begin{steps}[Conclusion of the Isomorphism]
Combining the three steps:
\begin{itemize}
    \item $\bar{f} \colon R/\ker(f) \to \operatorname{Im}(f)$ is a well-defined ring homomorphism (Step 1).
    \item $\bar{f}$ is surjective onto $\operatorname{Im}(f)$ (Step 2).
    \item $\bar{f}$ is injective, having trivial kernel (Step 3).
\end{itemize}
A bijective ring homomorphism is an isomorphism. Thus, $R/\ker(f) \cong \operatorname{Im}(f)$, which completes the proof of the First Isomorphism Theorem (\cref{thm:first-isomorphism-theorem-statement}).
\end{steps}
\end{content}

\end{proofWrapper}

\subsection{Recapitulation and the Anatomy of the Theorem}

\begin{speech}[01:06:32 - 01:07:37]
Those are the three steps: you define $\bar{f}$, you show it's surjective onto the image, and you show its kernel is zero.

The reason why the last one is confusing is you have to understand who is zero where. The kernel of $f$ is not zero; the kernel of $f$ is just whatever the kernel of $f$ is. But we quotient by the kernel of $f$, so the kernel of $\bar{f}$ is zero.

That's an admittedly confusing state of affairs, but you have to somehow come to terms with that.
\end{speech}

\begin{note}[Board Cleaning and Re-statement]
\inlinemetanote{The lecturer wipes the lower blackboard clean from 01:07:00 to 01:07:37 to write out the clean summary of the three objects and their respective zero elements.}
\end{note}

\begin{speech}[01:07:37 - 01:09:56]
I'll say it in words, because it's confusing. We start with $f \colon R \to S$. \inlinemetanote{writes $f \colon R \to S$ on the lower board} This has a kernel, $\ker(f) \trianglelefteq R$, \inlinemetanote{writes has kernel $\ker(f) \trianglelefteq R$} which could be anything --- maybe zero, maybe not, could be anything in the world.

Then, from this, we make $\bar{f} \colon R/\ker(f) \to \operatorname{Im}(f) \subseteq S$. \inlinemetanote{writes $\bar{f} \colon R/\ker(f) \to \operatorname{Im}(f) \subseteq S$} So, we study this. And the last step is that the kernel of $\bar{f}$ --- which is an ideal in $R/\ker(f)$ \inlinemetanote{writes $\ker(\bar{f}) \subseteq R/\ker(f)$} --- is the zero ideal in this quotient ring: $\ker(\bar{f}) = (0) \subseteq R/\ker(f)$. \inlinemetanote{writes $\ker(\bar{f}) = (0) \subseteq R/\ker(f)$}

The kernel of $f$ may not be zero, but when I consider the kernel of $\bar{f}$ in the quotient ring, it's zero!

And the way to think about that --- we'll discuss this more tomorrow --- is that you can write the kernel of $\bar{f}$ as kind of whatever was in the kernel before, but now quotiented out by $\ker(f)$, so that's just the zero group: \inlinemetanote{writes $\ker(\bar{f}) = \frac{\ker(f)}{\ker(f)} = 0$} $\frac{\ker(f)}{\ker(f)} = 0$. That's a nice way to think about it; I'll explain tomorrow how this kind of logic makes sense. Roughly speaking, there was a kernel, but when we quotient out, we quotient out by that kernel, so after we quotient out, we have no kernel left.

These three steps, taken together: first we defined $\bar{f}$, we showed it's surjective onto the image of $f$, and now we showed the kernel is zero, which means it's also injective. Kernel being zero implies that $\bar{f}$ is injective --- we started the lecture with that, I think. \inlinemetanote{writes $\implies \bar{f}$ is injective}

So, steps 1 to 3 show that $\bar{f} \colon R/\ker(f) \to \operatorname{Im}(f)$ \inlinemetanote{writes Steps 1--3: $\bar{f} \colon R/\ker(f) \to \operatorname{Im}(f)$} is well-defined, injective, and surjective, so $\bar{f}$ is an isomorphism of rings. \inlinemetanote{writes is isom of rings}

That's the end of the proof of the First Isomorphism Theorem.
\end{speech}

\begin{content}[Summary: Anatomy of the First Isomorphism Theorem]
The architecture of the First Isomorphism Theorem relies on distinguishing the three rings and their respective kernels:
\begin{center}
\begin{tabular}{l l l}
    \toprule
    \textbf{Mapping} & \textbf{Domain / Codomain} & \textbf{Kernel} \\
    \midrule
    $f$ & $R \longrightarrow S$ & $\ker(f) \trianglelefteq R$ (Arbitrary ideal of $R$) \\
    $\bar{f}$ & $R/\ker(f) \overset{\cong}{\longrightarrow} \operatorname{Im}(f)$ & $\ker(\bar{f}) = (\bar{0}) \trianglelefteq R/\ker(f)$ (The zero ideal in $R/\ker(f)$) \\
    \bottomrule
\end{tabular}
\end{center}

\begin{align}
f \colon R &\longrightarrow S && \text{with kernel } \ker(f) \trianglelefteq R, \label{eq:f-kernel-summary} \\
\bar{f} \colon R/\ker(f) &\overset{\sim}{\longrightarrow} \operatorname{Im}(f) \subseteq S && \text{with kernel } \ker(\bar{f}) = (0) \subseteq R/\ker(f). \label{eq:fbar-kernel-summary}
\end{align}
\end{content}

\begin{speech}[01:09:56 - 01:10:46]
There was a lot of chalk on the board, but if you really understand this, you'll see there's actually nothing happening in this proof: you just have to follow the definitions at every stage. It's similar to when you really understand that linear algebra result: the fact that the image is the space mod the kernel has almost nothing happening in that proof --- the definitions almost make it so, and that's what happens here also.

So, let's do an example.
\end{speech}

\section{Ideals Generated by Elements}

\begin{speech}[01:10:46 - 01:12:23]
I have some room here, so I can say: one of the things to think about rings, their ideals, and quotient rings is that we need to have some way of writing down ideals. Otherwise, we can't have too much fun with examples!

A standard way of writing an ideal is: let $x_1, \dots, x_n$ be elements of the ring $R$. \inlinemetanote{writes Let $x_1, \dots, x_n \in R$} They could also be infinite, but in this class typically it's finite.

If you have some list of elements in $R$, the standard notation is you put parentheses around them, and you define an ideal: \inlinemetanote{writes then $(x_1, \dots, x_n) \subseteq R$ is the ideal defined by $(x_1, \dots, x_n) = \{ \sum_{i=1}^n a_i x_i \mid a_i \in R \}$} this is the ideal defined by this equation, so I just have to tell you who is in this ideal.

This is equal to those elements that are of the form $a_1 x_1 + \dots + a_n x_n$. In words, you could say these are $R$-linear combinations of these elements: \inlinemetanote{writes $R$-linear combs of the $x_i$} $a_1 x_1 + \dots + a_n x_n$, where the $a_i$ are in $R$.

So, this is a particular set.
\end{speech}

\begin{content}[Definition of Ideals Generated by Elements]
Let $R$ be a commutative ring with identity $1$.

\begin{definition}[Ideal Generated by Finitely Many Elements]\label[definition]{def:ideal-generated-by-elements}
Let $x_1, \dots, x_n \in R$. The \newterm{ideal generated by} $x_1, \dots, x_n$, denoted by $(x_1, \dots, x_n)$, is the set of all $R$-linear combinations of the generators:
\begin{equation}\label{eq:finitely-generated-ideal-def}
(x_1, \dots, x_n) := \left\{ \sum_{i=1}^n a_i x_i \;\middle|\; a_1, \dots, a_n \in R \right\}.
\end{equation}
When $n = 1$, the ideal $(x) = \{ a x \mid a \in R \} = x R$ is called a \newterm{principal ideal}.
\end{definition}
\end{content}

\begin{speech}[01:12:23 - 01:13:39]
As I said, in words, you could say they're $R$-linear combinations of the $x_i$.

Another phrasing is that this is the \qt{ideal generated by the $x_i$}. \inlinemetanote{writes \qt{ideal generated by the $x_i$}} That's also language used: ideal generated by the $x_i$. Those are two different ways to express the same thing.

And it's just this set. You have to accept that this is an ideal, which means you have to show that zero is here --- of course zero is here if I take all these coefficients to be zero. You have to show that if something is here, then its negative is here: I just flip all the signs of all the $a_i$. You have to show that if something's here, when I add it to something else, it's still here: then I just add the coefficients of the linear combination. And finally, I have to show the multiplicative property: if something's here and I multiply it by $r$, then it's still there; but you just view that $r$ as a scalar, and it multiplies all these coefficients.

So, you have to check it --- actually, I just checked it for you!

This is a very nice way: if you take a ring, you can just pick any elements and you get an ideal. It's the ideal generated by that.

Are there questions about this? This is a really fundamental way of generating ideals.
\end{speech}

\begin{content}[Verification of the Ideal Axioms for \texorpdfstring{$(x_1, \dots, x_n)$}{(x1, ..., xn)}]
\begin{proposition}\label[proposition]{prop:finitely-generated-is-ideal}
For any elements $x_1, \dots, x_n \in R$, the set $(x_1, \dots, x_n)$ defined in \cref{def:ideal-generated-by-elements} is indeed an ideal of $R$. Moreover, it is the smallest ideal of $R$ containing the set $\{x_1, \dots, x_n\}$.
\end{proposition}

\begin{proof}
We check the two defining conditions of an ideal:
\begin{enumerate}[label=\textbf{(\arabic*)}]
    \item \textbf{Additive Subgroup:}
    \begin{itemize}
        \item Setting $a_1 = \dots = a_n = 0$ yields $0 = \sum_{i=1}^n 0 \cdot x_i \in (x_1, \dots, x_n)$.
        \item For any element $u = \sum_{i=1}^n a_i x_i$, its additive inverse is $-u = \sum_{i=1}^n (-a_i) x_i \in (x_1, \dots, x_n)$.
        \item If $u = \sum_{i=1}^n a_i x_i$ and $v = \sum_{i=1}^n b_i x_i$, their sum is:
        \[
        u + v = \sum_{i=1}^n (a_i + b_i) x_i \in (x_1, \dots, x_n).
        \]
    \end{itemize}
    \item \textbf{Multiplicative Absorption:}
    For any ring element $r \in R$ and $u = \sum_{i=1}^n a_i x_i \in (x_1, \dots, x_n)$:
    \[
    r \cdot u = r \sum_{i=1}^n a_i x_i = \sum_{i=1}^n (r a_i) x_i.
    \]
    Since $R$ is a ring, each coefficient $r a_i \in R$, so $r \cdot u \in (x_1, \dots, x_n)$.
\end{enumerate}
\end{proof}
\end{content}

\subsection{Quotients of Univariate Polynomial Rings}

\begin{speech}[01:13:39 - 01:14:56]
I have a few minutes, so let's do some examples.

Now you have a very flexible way of generating a lot of rings. And the rings that you have the power to generate already have all sorts of interesting, different behavior. Part of the course is exploring some of that structure --- by far not all of it!
\end{speech}

\begin{note}[Board Cleaning and Transition to Examples]
\inlinemetanote{The lecturer washes the left and middle blackboards with a sponge to prepare for concrete polynomial examples.}
\end{note}

\begin{speech}[01:14:56 - 01:16:26]
For example, I can take my ring $R$ to be polynomials in one variable: $R = \mathbb{C}[x]$. \inlinemetanote{writes $R = \mathbb{C}[x]$} You should be totally happy with this ring: it's polynomials with complex coefficients.

And I can take an element here, $f$, to be $x^2 + 1$. That's a polynomial. \inlinemetanote{writes $f = x^2 + 1 \in R$}

And then I can take the ideal --- I'm using the notation from before: if you have one element, that's an element; if I put parentheses around it, this is the ideal: $(f) = (x^2 + 1) \subseteq R$. \inlinemetanote{writes $(f) = (x^2+1) \subseteq R$ ideal}

Very concretely, using that definition, this is just the ideal of $f$ times a polynomial, where the polynomial is in the ring: \inlinemetanote{writes $(f) = \{ f \cdot p \mid p \in R \}$} it's just one linear combination. All this is, is polynomials that have an $x^2 + 1$ factor (i.e., multiples of $x^2 + 1$).

What we've done today is: we have a new ring now, $R$ modulo this ideal. \inlinemetanote{writes $R/(f) = \mathbb{C}[x]/(x^2+1)$} This is the quotient ring, $\mathbb{C}[x]/(x^2 + 1)$.

Are you happy to think about that? We take this polynomial ring, we take this ideal, and we have this quotient ring. Do you accept that all the symbols on this board have been defined?

So, I can ask you now any question about this ring. For example: is this an integral domain? \inlinemetanote{writes Is this an integral domain?}

What's your opinion about that?
\end{speech}

\begin{student}[Student Answer]
No, because $x-i$ times $x+i$... $x-i$ is not zero and $x+i$ is not zero.
\end{student}

\begin{hidden}
% timestamp: 01:16:26
% topic: Example: R = C[x] and the quotient C[x]/(x^2+1); student identifies zero divisors.
% board_state: def:ideal-generated-by-elements, prop:finitely-generated-is-ideal
% next_goal: Formalize the zero divisor argument in C[x]/(x^2+1); contrast with Q[x]/(x^2+1).
% open_loops: none
\end{hidden}

\begin{speech}[01:16:26 - 01:18:39]
Yeah, the answer is no! \inlinemetanote{writes No!}

And the reason is: I can take $\overline{x + i} \cdot \overline{x - i}$, \inlinemetanote{writes $\overline{x+i} \cdot \overline{x-i}$} because $x + i$ and $x - i$ are in the polynomial ring, so their classes are in the quotient ring. And by definition, this is equal to $\overline{(x + i)(x - i)} = \overline{x^2 + 1}$, \inlinemetanote{writes $= \overline{(x+i)(x-i)} = \overline{x^2+1}$} and this is zero: $\bar{0}$, right? \inlinemetanote{writes $= \bar{0}$}

This is the beginning of your argument. That's a totally solid start!

To prove that it's not an integral domain, you have to find two non-zero things that multiply to zero. You found two things that multiply to zero. Why are these non-zero?
\end{speech}

\begin{student}[Student Explanation]
Because the difference to zero... if you multiply $x^2 + 1$ by any polynomial, the degree is at least $2$.
\end{student}

\begin{speech}[01:18:39 - 01:18:56]
That's the correct argument! We have to show that $\overline{x + i} \neq \bar{0}$ and $\overline{x - i} \neq \bar{0}$. \inlinemetanote{writes But $\overline{x+i} \neq \bar{0}$ and $\overline{x-i} \neq \bar{0}$} If we have these three statements --- the product equals zero, and they're both not equal to zero ---, then it's not an integral domain. You agree with that?

And he gave the correct argument for that: because if $x + i$ were equivalent to zero, this would mean that $x + i$ is in this ideal generated by $x^2 + 1$. \inlinemetanote{writes $\implies x+i \in (x^2+1)$}

But the ideal generated by $x^2 + 1$ has $(x^2 + 1)$ times polynomials in it. And you can never lower the degree to a linear polynomial by taking a quadratic polynomial and multiplying it by another polynomial. So, the full argument says: the answer is no.
\end{speech}

\begin{content}[Zero Divisors in the Quotient Ring \texorpdfstring{$\mathbb{C}[x]/(x^2+1)$}{C[x]/(x^2+1)}]
\begin{example}[Non-Domain Quotient Ring]\label[example]{ex:cx-quotient-not-domain}
Consider the polynomial ring $R = \mathbb{C}[x]$ and the principal ideal generated by $f(x) = x^2 + 1$:
\[
I = (x^2 + 1) = \{ (x^2 + 1) p(x) \mid p(x) \in \mathbb{C}[x] \}.
\]
The quotient ring $\mathbb{C}[x]/(x^2 + 1)$ is \textbf{not an integral domain}.
\end{example}

\begin{proof}
In $\mathbb{C}[x]$, the polynomial $x^2 + 1$ factors into linear factors:
\[
x^2 + 1 = (x + i)(x - i).
\]
In the quotient ring $\mathbb{C}[x]/(x^2 + 1)$, consider the cosets $\overline{x + i}$ and $\overline{x - i}$:
\begin{enumerate}[label=\textbf{(\arabic*)}]
    \item \textbf{Their Product is Zero:}
    \[
    \overline{x + i} \cdot \overline{x - i} = \overline{(x + i)(x - i)} = \overline{x^2 + 1} = \bar{0} \in \mathbb{C}[x]/(x^2 + 1).
    \]
    \item \textbf{Neither Factor is Zero:}
    Suppose $\overline{x + i} = \bar{0}$. Then $x + i \in (x^2 + 1)$, which means there exists $p(x) \in \mathbb{C}[x]$ such that:
    \[
    x + i = (x^2 + 1) p(x).
    \]
    If $p(x) = 0$, the right-hand side is $0 \neq x + i$. If $p(x) \neq 0$, comparing degrees gives:
    \[
    1 = \deg(x + i) = \deg(x^2 + 1) + \deg(p) = 2 + \deg(p) \ge 2,
    \]
    a contradiction. Thus, $\overline{x + i} \neq \bar{0}$. By identical degree considerations, $\overline{x - i} \neq \bar{0}$.
\end{enumerate}
Since $\overline{x + i}$ and $\overline{x - i}$ are non-zero elements whose product is zero, they are zero divisors. Hence, $\mathbb{C}[x]/(x^2 + 1)$ is not an integral domain.
\end{proof}
\end{content}

\subsection{The Quotient \texorpdfstring{$\mathbb{Q}[x]/(x^2+1)$}{Q[x]/(x^2+1)} and the Gaussian Rationals}

\begin{speech}[01:18:56 - 01:20:56]
What about this one? This is kind of cutting closer to what this class is about, actually: $\mathbb{Q}[x]/(x^2+1)$. \inlinemetanote{writes $\mathbb{Q}[x]/(x^2+1)$}

Is this an integral domain? \inlinemetanote{writes Is this an integral domain?}

His argument you cannot use here, because $i$ is not a rational number, right? You agree you can't use the same argument.

So, I can ask the same question, though: what do you think the answer to that is?

This is really a question that the last part of the class is about, but I'm going to give you a proof of this. The answer is: this is an integral domain! \inlinemetanote{writes is an int. domain}

Actually, more than that! \inlinemetanote{laughs} It's a field! \inlinemetanote{writes Field}

It doesn't look like it immediately, but I will try to give you a proof in the next few minutes. There are a lot of ways to prove it, and a certain part of the class is about such things called field extensions. If you take the successor course in the spring, that's entirely about this kind of stuff: Galois theory, etc. But it's a very good example to think about.

So far, the only thing you're supposed to understand is that the argument proposed for the complex numbers does not work here. But let's try to find an argument that does work, that's related to this lecture.

In some sense, that question is kind of like searching for the complex numbers: because you know the roots of $x^2 + 1$, built into that question is this quest for the complex numbers.
\end{speech}

\begin{content}[The Rational Quotient \texorpdfstring{$\mathbb{Q}[x]/(x^2+1)$}{Q[x]/(x^2+1)}]
\begin{question}
Is the quotient ring $\mathbb{Q}[x]/(x^2 + 1)$ an integral domain?
\end{question}

\begin{answer}
\textbf{Yes! In fact, it is a field.}
Notice that the factorization $x^2 + 1 = (x + i)(x - i)$ used over $\mathbb{C}$ is unavailable in $\mathbb{Q}[x]$, because $i \notin \mathbb{Q}$. In fact, $x^2 + 1$ is \newterm{irreducible} over $\mathbb{Q}$. We will prove that $\mathbb{Q}[x]/(x^2 + 1)$ is isomorphic to the field of Gaussian rationals $\mathbb{Q}(i)$.
\end{answer}
\end{content}

\subsection{Construction of the Evaluation Homomorphism at \texorpdfstring{$i$}{i}}

\begin{speech}[01:20:56 - 01:22:36]
But we know the complex numbers, so we write that down. \inlinemetanote{writes $\mathbb{C}$} And I take $\mathbb{Q}$: \inlinemetanote{writes $\mathbb{Q} \hookrightarrow \mathbb{C}$} $\mathbb{Q}$ embeds in the complex numbers; there's no doubt about that.

And now I look at the polynomial ring: $\mathbb{Q}[x]$. \inlinemetanote{writes $\mathbb{Q}[x]$} And we proved that if I choose any value for $x$, I get a homomorphism from the polynomial ring. Do you remember that?

So, I'm going to choose a homomorphism $f$ that on constants just acts as the inclusion, but $f(x) = i$. \inlinemetanote{writes $f \colon \mathbb{Q}[x] \to \mathbb{C}$, $f(x) = i$} I get to do it! Because I told you: for the polynomial ring, if I know how the base ring maps to any other ring, I can always extend the map sending $x$ to anything I want. So, I could just do this --- it's my right from previous lectures.

In other words, I define a homomorphism $f$ from polynomials to the complex numbers, \inlinemetanote{writes $f \colon \mathbb{Q}[x] \to \mathbb{C}$} where $f(p(x)) = p(i)$, \inlinemetanote{writes $f(p(x)) = p(i)$} okay? Do you accept this?

Even if you forgot that theorem, you can just check that this is a legal definition: given any polynomial, evaluating at $i$ gives a complex number, and it's a ring homomorphism. \inlinemetanote{writes Ring hom} So, you can just prove that from scratch in your head because I gave you the formula, or you can remember it has to be true because you can send $x$ anywhere.
\end{speech}

\begin{hidden}
% timestamp: 01:22:36
% topic: Defining the evaluation homomorphism f: Q[x] -> C, p(x) -> p(i).
% board_state: def:evaluation-homomorphism-i
% next_goal: Compute Im(f) = Q(i) and apply the First Isomorphism Theorem.
% open_loops: none
\end{hidden}

\begin{content}[The Evaluation Homomorphism at \texorpdfstring{$i$}{i}]
Let $\iota \colon \mathbb{Q} \hookrightarrow \mathbb{C}$ denote the canonical inclusion of fields. By the universal property of polynomial rings, specifying the image of the indeterminate $x$ uniquely determines a ring homomorphism extending $\iota$.

\begin{center}
\begin{tikzpicture}[scale=1.2, >=Stealth]
    % \begin{hidden}
    % - Canvas: scale 1.2, commutative triangle with Q at (0,1.5), C at (3.5,1.5), Q[x] at (0,0)
    % - Elements: inclusion Q into C, canonical inclusion Q into Q[x], evaluation homomorphism f from Q[x] to C
    % - Labels: iota on top, inc on left, f: p(x) -> p(i) on hypotenuse
    % \end{hidden}
    \node (Q) at (0, 1.5) {$\mathbb{Q}$};
    \node (C) at (3.5, 1.5) {$\mathbb{C}$};
    \node (Qx) at (0, 0) {$\mathbb{Q}[x]$};
    
    \draw[right hook->, thick] (Q) -- node[above] {$\iota$} (C);
    \draw[right hook->, thick] (Q) -- node[left] {$\mathrm{inc}$} (Qx);
    \draw[->, thick, MidnightBlue] (Qx) -- node[below right] {$f \colon p(x) \mapsto p(i)$} (C);
\end{tikzpicture}
\end{center}

\begin{definition}[Evaluation Homomorphism at $i$]\label[definition]{def:evaluation-homomorphism-i}
We define the \newterm{evaluation homomorphism} at the imaginary unit $i \in \mathbb{C}$ by:
\begin{equation}\label{eq:eval-homomorphism-def}
f \colon \mathbb{Q}[x] \longrightarrow \mathbb{C}, \qquad f(p(x)) := p(i).
\end{equation}
Explicitly, if $p(x) = \sum_{k=0}^m a_k x^k \in \mathbb{Q}[x]$, then $f(p(x)) = \sum_{k=0}^m a_k i^k \in \mathbb{C}$.
\end{definition}

\begin{steps}[Homomorphism Property Verification]
For any $p(x), q(x) \in \mathbb{Q}[x]$:
\begin{align*}
f(p(x) + q(x)) &= (p + q)(i) = p(i) + q(i) = f(p(x)) + f(q(x)), \\
f(p(x) \cdot q(x)) &= (p \cdot q)(i) = p(i) \cdot q(i) = f(p(x)) \cdot f(q(x)), \\
f(1) &= 1.
\end{align*}
Thus, $f$ is a well-defined homomorphism of commutative rings with identity.
\end{steps}
\end{content}

\subsection{Image of the Evaluation Homomorphism and the First Isomorphism Theorem}

\begin{speech}[01:22:36 - 01:23:54]
What is the image of this? \inlinemetanote{writes What is $\operatorname{Im}(f) = ?$ on the middle board} What is the image? Do you know what that is? Yeah?
\end{speech}

\begin{student}[Student Answer]
I mean, if we consider it like polynomials of degree $1$ with a constant, we get the Gaussian rationals?
\end{student}

\begin{speech}[01:23:54 - 01:24:51]
Yeah, and you don't get any more than that! Because if you take a polynomial, all the quadratic terms give you $-1$, and all you get is $i$. So, the answer is: the image of this is what we'd call the Gaussian integers (i.e., Gaussian rationals $\mathbb{Q}(i)$). \inlinemetanote{writes $\mathbb{Q}(i)$}

If you take a polynomial with rational coefficients and you evaluate it at $i$, all the even terms just give you rational numbers, and all the odd terms give you $i$'s. And the coefficients are all rational numbers, and you can get any one you want here.

So, now we use our First Isomorphism Theorem! \inlinemetanote{writes First Isom. Theorem} That says that $\mathbb{Q}[x]$ modulo the kernel of $f$ is isomorphic to the image of $f$, which is the Gaussian integers (i.e., Gaussian rationals $\mathbb{Q}(i)$) --- which, by the way, is a field! It's a domain and it's a field: \inlinemetanote{writes $\mathbb{Q}[x]/\ker(f) \cong \operatorname{Im}(f) = \mathbb{Q}(i)$} $\mathbb{Q}[x]/\ker(f) \cong \operatorname{Im}(f) = \mathbb{Q}(i)$.
\end{speech}

\begin{content}[Image of $f$ and Application of the First Isomorphism Theorem]
\begin{proposition}[Image of the Evaluation Map]\label[proposition]{prop:image-eval-map}
The image of the evaluation homomorphism $f \colon \mathbb{Q}[x] \to \mathbb{C}$ is the field of Gaussian rationals:
\begin{equation}\label{eq:image-gaussian-rationals}
\operatorname{Im}(f) = \mathbb{Q}(i) := \{ a + b i \mid a, b \in \mathbb{Q} \} \subset \mathbb{C}.
\end{equation}
\end{proposition}

\begin{proof}
Let $p(x) = a_n x^n + \dots + a_1 x + a_0 \in \mathbb{Q}[x]$. Since $i^2 = -1$, every even power $i^{2k} = (-1)^k \in \mathbb{Q}$ and every odd power $i^{2k+1} = (-1)^k i$. Collecting the even and odd terms yields:
\[
p(i) = \left( \sum_{k \ge 0} a_{2k} (-1)^k \right) + \left( \sum_{k \ge 0} a_{2k+1} (-1)^k \right) i = a + b i,
\]
where $a, b \in \mathbb{Q}$. Conversely, for any $a + b i \in \mathbb{Q}(i)$, the linear polynomial $p(x) = b x + a \in \mathbb{Q}[x]$ satisfies $f(p(x)) = a + b i$. Thus, $\operatorname{Im}(f) = \mathbb{Q}(i)$.
\end{proof}

\begin{theorem}[First Isomorphism Theorem Application]\label[theorem]{thm:first-isom-applied-qi}
Applying the First Isomorphism Theorem (\cref{thm:first-isomorphism-theorem-statement}) to $f \colon \mathbb{Q}[x] \to \mathbb{C}$:
\begin{equation}\label{eq:isom-qi-statement}
\frac{\mathbb{Q}[x]}{\ker(f)} \cong \operatorname{Im}(f) = \mathbb{Q}(i).
\end{equation}
Since $\mathbb{Q}(i)$ is a subfield of $\mathbb{C}$, the quotient $\mathbb{Q}[x]/\ker(f)$ is an integral domain and a field.
\end{theorem}
\end{content}

\begin{aiNote}[Terminology Note: Gaussian Rationals vs.\ Gaussian Integers]
% [PERSISTENT]
In the lecture audio, the professor refers to $\mathbb{Q}(i) = \{ a + b i \mid a, b \in \mathbb{Q} \}$ as the \qt{Gaussian integers}. In standard algebraic terminology, the Gaussian integers are $\mathbb{Z}[i] = \{ a + b i \mid a, b \in \mathbb{Z} \}$, while $\mathbb{Q}(i) = \operatorname{Frac}(\mathbb{Z}[i])$ is the field of \newterm{Gaussian rationals}. The blackboard notation $\mathbb{Q}(i)$ is preserved verbatim.
\end{aiNote}

\subsection{Computing the Kernel via Polynomial Long Division}

\begin{speech}[01:24:51 - 01:27:11]
Now it would be nice to know what this kernel is: what is the kernel of $f$? \inlinemetanote{writes What is the ker of $f$? on the lower right board}

Well, one thing that's definitely in the kernel is this polynomial: $x^2 + 1$ is definitely in the kernel, \inlinemetanote{writes $x^2 + 1 \in \ker(f)$} because $f(x^2 + 1)$, by definition, is equal to $i^2 + 1 = 0$, right? \inlinemetanote{writes $f(x^2+1) = i^2 + 1 = 0$} So, whatever the kernel is, it definitely contains this polynomial.

Once this polynomial is in the kernel, then the whole ideal generated by this polynomial is also in the kernel: \inlinemetanote{writes $(x^2+1) \subseteq \ker(f)$} $(x^2 + 1) \subseteq \ker(f)$, because if I take this polynomial and multiply it by anything else, it will still be in the kernel.

The question now is: is there anything else? Is the ideal generated by this polynomial equal to the kernel of $f$? \inlinemetanote{writes Is $(x^2+1) = \ker(f)$?}

If the answer is yes, then I've answered that question! Spoiler: it's going to be yes. If it is yes, then I will have answered this question, because I've computed the kernel here and shown it's isomorphic to the field of Gaussian integers (i.e., rationals $\mathbb{Q}(i)$). That is a way of writing the field of Gaussian rationals in a slightly hidden form.

So, all I have to do is prove that $(x^2 + 1)$ is the kernel.
\end{speech}

\begin{content}[Containment of the Principal Ideal $(x^2+1)$ in the Kernel]
\begin{lemma}[Ideal Containment in the Kernel]\label[lemma]{lem:ideal-contained-in-kernel}
The principal ideal $(x^2 + 1)$ is contained in $\ker(f)$:
\begin{equation}\label{eq:ideal-contained-ker}
(x^2 + 1) \subseteq \ker(f).
\end{equation}
\end{lemma}

\begin{proof}
Evaluating $f$ on the generator $x^2 + 1$:
\[
f(x^2 + 1) = i^2 + 1 = -1 + 1 = 0 \implies x^2 + 1 \in \ker(f).
\]
Since $\ker(f)$ is an ideal of $\mathbb{Q}[x]$ (\cref{prop:kernel-is-ideal}), it must absorb multiplication by any polynomial $q(x) \in \mathbb{Q}[x]$:
\[
\forall q(x) \in \mathbb{Q}[x] \implies q(x)(x^2 + 1) \in \ker(f).
\]
Hence, $(x^2 + 1) \subseteq \ker(f)$.
\end{proof}
\end{content}

\begin{speech}[01:27:11 - 01:28:36]
We have $f \colon \mathbb{Q}[x] \to \mathbb{C}$, where $f$ eats a polynomial and gives back a complex number, $p(i)$. A polynomial is in the kernel of $f$ if and only if $p(i) = 0$. \inlinemetanote{writes $p \in \ker(f) \iff p(i) = 0$} In real life, you say $i$ is a root of $p$, right? \inlinemetanote{writes \qt{$i$ is a root of $p$}} So, the polynomials in the kernel are those where $i$ is a root.

We have to see how that's related to $x^2 + 1$. We could do the following: we try to divide $p(x)$ by $x^2 + 1$, \inlinemetanote{writes $x^2+1 \mid p(x)$ on the lower board} and everything is happening in this ring of polynomials $\mathbb{Q}[x]$. \inlinemetanote{writes in $\mathbb{Q}[x]$}

This is what's called polynomial long division, which you should have learned in school. If you have a polynomial, you can start dividing it: it will be equal to some polynomial $q(x)$ times $x^2 + 1$ plus a remainder term. \inlinemetanote{writes $p(x) = q(x)(x^2+1) + r(x)$}

Do you remember polynomial long division? You just keep taking out the highest powers all the way down. If you do polynomial long division, you start dividing and you'll eventually get some polynomial times $(x^2 + 1)$ plus a remainder term. And since this is a quadratic polynomial, this remainder term is a linear polynomial: \inlinemetanote{writes linear poly $r_1 x + r_2$, $r_1, r_2 \in \mathbb{Q}$} some $r_1 x + r_2$, where $r_1, r_2 \in \mathbb{Q}$. Of course, these could be zero.
\end{speech}

\begin{content}[Polynomial Division by \texorpdfstring{$x^2+1$}{x^2+1}]
\begin{lemma}[{Division Algorithm in $\mathbb{Q}[x]$}]\label[lemma]{lem:poly-division-algorithm}
For any polynomial $p(x) \in \mathbb{Q}[x]$, there exist unique polynomials $q(x), r(x) \in \mathbb{Q}[x]$ such that:
\begin{equation}\label{eq:poly-division-x2-plus-1}
p(x) = q(x)(x^2 + 1) + r(x), \quad \text{with } \deg(r) < 2.
\end{equation}
In particular, the remainder $r(x)$ is a linear polynomial with rational coefficients:
\begin{equation}\label{eq:remainder-linear-form}
r(x) = r_1 x + r_2, \qquad r_1, r_2 \in \mathbb{Q}.
\end{equation}
\end{lemma}
\end{content}

\begin{speech}[01:28:36 - 01:29:21]
That's polynomial long division; that's something you learned in school, and you should remember it now.

Now, I just stick $i$ into this equation: \inlinemetanote{writes $p(i) = q(i)(i^2+1) + r(i)$} I put $i$'s everywhere: $p(i) = q(i)(i^2 + 1) + r(i)$.

And now, the assumption is $p(i) = 0$, so the left-hand side is zero. This first term is also zero because $i^2 + 1 = 0$. That means the conclusion is: $r(i) = 0$. \inlinemetanote{writes $0 = q(i) \cdot 0 + r(i) \implies r(i) = 0$}

How could that be? The remainder $r$ is a linear polynomial. How could $r(i)$ be zero? There's only one way it could be zero, which is what? Yeah?
\end{speech}

\begin{student}[Student Answer]
Both coefficients are zero.
\end{student}

\begin{speech}[01:29:21 - 01:30:14]
Yeah, it's the only way! If this were a non-zero polynomial, since these coefficients are rational numbers, the only way this is zero when you stick in $i$ is if this is the zero polynomial, because $i$ is not a rational number.

So, this implies that $r$ is the zero polynomial! \inlinemetanote{writes $\implies r = 0$} So, actually it's not here, and $p(x)$ is in the ideal! \inlinemetanote{points to $p(x) = q(x)(x^2+1)$} And that finishes the argument, actually.

The conclusion is that $p(i) = 0$ implies $p(x) \in (x^2 + 1)$. \inlinemetanote{writes $p(i) = 0 \implies p(x) \in (x^2+1)$}

And the final conclusion, using the First Isomorphism Theorem, is that $\mathbb{Q}[x]/(x^2 + 1)$ \inlinemetanote{writes $\mathbb{Q}[x]/(x^2+1)$} is isomorphic \inlinemetanote{writes $\cong \mathbb{Q}(i)$} to this ring of Gaussian integers (i.e., Gaussian rationals $\mathbb{Q}(i)$).

Okay, made it in time! That's it. \inlinemetanote{lecturer finishes and packs up}
\end{speech}

\begin{content}[Proof that \texorpdfstring{$\ker(f) = (x^2+1)$}{ker(f) = (x^2+1)} and Final Isomorphism]
\begin{theorem}[Equality of Kernel and Principal Ideal]\label[theorem]{thm:kernel-equals-ideal}
The kernel of the evaluation homomorphism $f \colon \mathbb{Q}[x] \to \mathbb{C}$ at $i$ is precisely the principal ideal generated by $x^2 + 1$:
\begin{equation}\label{eq:kernel-equals-ideal}
\ker(f) = (x^2 + 1).
\end{equation}
\end{theorem}

\begin{proof}
We established in \cref{lem:ideal-contained-in-kernel} that $(x^2 + 1) \subseteq \ker(f)$. We now prove the reverse inclusion $\ker(f) \subseteq (x^2 + 1)$.

Let $p(x) \in \ker(f)$. By definition, $p(i) = 0$. By \cref{lem:poly-division-algorithm}, divide $p(x)$ by $x^2 + 1$:
\[
p(x) = q(x)(x^2 + 1) + r(x), \quad \text{where } r(x) = r_1 x + r_2 \text{ with } r_1, r_2 \in \mathbb{Q}.
\]
Evaluating both sides at $x = i$:
\[
p(i) = q(i)(i^2 + 1) + r(i).
\]
Since $p(i) = 0$ and $i^2 + 1 = -1 + 1 = 0$, we obtain:
\[
0 = q(i) \cdot 0 + r(i) \implies r(i) = 0.
\]
Substituting the linear form of $r(x)$:
\[
r(i) = r_1 i + r_2 = 0 \implies r_2 + r_1 i = 0.
\]
Since $r_1, r_2 \in \mathbb{Q}$ and $i \notin \mathbb{Q}$ (so $1$ and $i$ are linearly independent over $\mathbb{Q}$), this forces:
\[
r_1 = 0 \quad \text{and} \quad r_2 = 0 \implies r(x) = 0.
\]
Therefore, the remainder vanishes identically, and:
\[
p(x) = q(x)(x^2 + 1) \in (x^2 + 1).
\]
Thus, $\ker(f) \subseteq (x^2 + 1)$, which completes the proof that $\ker(f) = (x^2 + 1)$.
\end{proof}

\begin{corollary}[Realization of the Gaussian Rationals as a Quotient Ring]\label[corollary]{cor:gaussian-rationals-quotient}
Combining \cref{thm:first-isom-applied-qi} and \cref{thm:kernel-equals-ideal}, there is a canonical isomorphism of fields:
\begin{equation}\label{eq:gaussian-rationals-quotient-isom}
\frac{\mathbb{Q}[x]}{(x^2 + 1)} \overset{\cong}{\longrightarrow} \mathbb{Q}(i), \qquad \overline{p(x)} \mapsto p(i).
\end{equation}
In particular, $\mathbb{Q}[x]/(x^2 + 1)$ is an integral domain and a field.
\end{corollary}
\end{content}

\begin{overview}[Summary of Key Concepts]
\begin{itemize}
    \item \textbf{Polynomial and Power Series Rings:} The rings $R[x]$ and $R[[x]]$ preserve the integral domain property of $R$, relying on the highest-degree terms and lowest-degree orders, respectively.
    \item \textbf{Laurent Series and Rational Functions:} The field of rational functions $R(x)$ is $\operatorname{Frac}(R[x])$, while the field of formal Laurent series $K((x))$ (finite negative powers) is naturally isomorphic to $\operatorname{Frac}(K[[x]])$.
    \item \textbf{Ideals and Well-Defined Quotients:} An ideal $I \trianglelefteq R$ requires additive closure and multiplicative absorption ($R \cdot I \subseteq I$). In constructing the quotient ring $R/I$, the additive subgroup condition guarantees well-defined addition, while the absorption property is strictly required to ensure well-defined coset multiplication.
    \item \textbf{The First Isomorphism Theorem for Rings:} Every ring homomorphism $f \colon R \to S$ induces a canonical ring isomorphism $\bar{f} \colon R/\ker(f) \overset{\cong}{\to} \operatorname{Im}(f)$ via $\bar{x} \mapsto f(x)$.
    \item \textbf{Polynomial Quotients:} Modding out by $x^2+1$ over $\mathbb{C}$ produces zero divisors via $(x+i)(x-i) \equiv 0$, whereas over $\mathbb{Q}$, polynomial irreducibility and evaluation at $i$ yield the field of Gaussian rationals $\mathbb{Q}[x]/(x^2+1) \cong \mathbb{Q}(i)$.
\end{itemize}
\end{overview}